The CKKS scheme supports SIMD approximate arithmetic on complex numbers, including real numbers as a special case. It is useful for numerical workloads such as machine learning. BFV, BGV, and TFHE have exact finite plaintext spaces, but can also implement fixed-point or floating-point arithmetic through suitable encodings or circuits; the distinction is their native arithmetic and error semantics.~\cite{ckks-original}

For a power-of-two $n\geq 2$, distinguish the real algebra $\mathcal{R}_{\langle n\rangle}=\mathbb{R}[X]/(X^n+1)$, the integer ring $R=\mathbb{Z}[X]/(X^n+1)$, and the ciphertext ring $R_q=R/qR$. The inverse canonical embedding maps $n/2$ complex slots to a polynomial $M\in\mathcal{R}_{\langle n\rangle}$. Encoding then scales its coefficients by $\Delta>0$ and rounds them to obtain an integer polynomial in $R$, which is encrypted modulo $q$.


$ $

\noindent Overall, CKKS's encryption procedure is as follows:

\begin{enumerate}
\item \underline{\textsf{Encoding\textsubscript{1}}:} Encode the targeted input complex number as a real number
\item \underline{\textsf{Encoding\textsubscript{2}}:} Encode the real number as an integer
\item \underline{\textsf{Encryption}:} Encrypt the integer using RLWE
\end{enumerate}
The encrypted RLWE ciphertext supports homomorphic addition and multiplication.

$ $

\noindent At the end of all homomorphic operations, CKKS's decryption procedure is as follows:

\begin{enumerate}
\item \underline{\textsf{Decryption}:} Decrypt the RLWE ciphertext into a plaintext integer
\item \underline{\textsf{Decoding\textsubscript{1}}:} Decode the integer to a real number
\item \underline{\textsf{Decoding\textsubscript{2}}:} Decode the real number to a complex number
\end{enumerate}

$ $

CKKS generally introduces encoding error when the scaled inverse embedding is rounded to the integer lattice. Encryption noise and subsequent arithmetic, key-switching, rescaling, and approximation errors also contribute to the decoded error. This is not specifically an error from square roots, and some inputs can be encoded exactly. Unlike exact plaintext decryption in BFV, general CKKS decryption retains the numerical error.

CKKS and BFV share ring arithmetic and key-switching techniques, but their encoding, decoding, scale management, multiplication, and bootstrapping must respect different plaintext semantics. Both use rings: $\mathbb{C}$ is itself a field, hence a ring. In CKKS, the relevant distinction is approximate arithmetic in characteristic zero instead of exact arithmetic modulo a plaintext modulus.



\begin{tcolorbox}[
    title = \textbf{Required Background},    % box title
    colback = white,    % light background; tweak to taste
    colframe = black,  % frame colour
    boxrule = 0.8pt,     % line thickness
    left = 1mm, right = 1mm, top = 1mm, bottom = 1mm % inner padding
]

\begin{itemize}
\item \autoref{sec:modulo}: \nameref{sec:modulo}
\item \autoref{sec:group}: \nameref{sec:group}
\item \autoref{sec:field}: \nameref{sec:field}
\item \autoref{sec:order}: \nameref{sec:order}
\item \autoref{sec:polynomial-ring}: \nameref{sec:polynomial-ring}
\item \autoref{sec:decomp}: \nameref{sec:decomp}
\item \autoref{sec:roots}: \nameref{sec:roots}
\item \autoref{sec:cyclotomic}: \nameref{sec:cyclotomic}
\item \autoref{sec:matrix}: \nameref{sec:matrix}
\item \autoref{sec:euler}: \nameref{sec:euler}
\item \autoref{sec:modulus-rescaling}: \nameref{sec:modulus-rescaling}
\item \autoref{sec:chinese-remainder}: \nameref{sec:chinese-remainder}
\item \autoref{sec:taylor-series}: \nameref{sec:taylor-series}
\item \autoref{sec:polynomial-interpolation}: \nameref{sec:polynomial-interpolation}
\item \autoref{sec:ntt}: \nameref{sec:ntt}
\item \autoref{sec:lattice}: \nameref{sec:lattice}
\item \autoref{sec:rlwe}: \nameref{sec:rlwe}
\item \autoref{sec:glwe}: \nameref{sec:glwe}
\item \autoref{sec:glev}: \nameref{sec:glev}
\item \autoref{sec:glwe-add-cipher}: \nameref{sec:glwe-add-cipher}
\item \autoref{sec:glwe-add-plain}: \nameref{sec:glwe-add-plain}
\item \autoref{sec:glwe-mult-plain}: \nameref{sec:glwe-mult-plain}
\item \autoref{subsec:modulus-switch-rlwe}: \nameref{subsec:modulus-switch-rlwe}
\item \autoref{sec:glwe-key-switching}: \nameref{sec:glwe-key-switching}
\item \autoref{sec:bfv}: \nameref{sec:bfv}
\end{itemize}
\end{tcolorbox}

\clearpage


\subsection{Encoding and Decoding}
\label{subsec:ckks-encoding-decoding}


CKKS's encoding and decoding is fundamentally very similar to BFV's batch encoding scheme. BFV designs its batch encoding scheme (Summary~\ref*{subsec:bfv-enc-dec} in \autoref{subsec:bfv-enc-dec}) based on the updated $\hathat W$ and $\hathat W^*$ matrices (Summary~\ref*{subsubsec:bfv-rotation-summary} in \autoref{subsec:bfv-rotation}). That is, BFV decodes a polynomial into an input slot vector by evaluating the polynomial at each root of $X^n+1$, which is the primitive $(\mu=2n)$-th root of unity (i.e., $\vec{v} = \hathat W^* \cdot \vec{m}$), and encodes an input slot vector into a polynomial by inversing this operation (i.e., $\vec{m} = n^{-1} \cdot \hathat W \cdot I_n^R \cdot \vec{v}$). This encoding and decoding scheme is designed based on Summary~\ref*{subsec:poly-vector-transformation-complex} (\autoref{subsec:poly-vector-transformation-complex}) which designs the isomorphic mapping between $n$-slot vectors in a ring (finite field) and $(n-1)$-degree (or lesser degree) polynomials as follows:

$\sigma: f(x) \in \mathbb{Z}_t[X] / F(X) \text{ } \longrightarrow \text{ } (f(\omega^1), f(\omega^3)), \cdots, f(\omega^{2n-1})) \in \mathbb{Z}_t^n$

, where $\omega = g^{\frac{t - 1}{2n}}$ is a root of (i.e., primitive $(\mu=2n)$-th root of unity) of the $(\mu=2n)$-th cyclotomic polynomial $X^n + 1$ defined over a prime modulo $t$ ring.

$ $

CKKS uses the same evaluation/interpolation principle over complex roots of unity. Complex numbers also form a ring; here the full evaluation vector lies in the real algebra $\mathbb{\hat C}^n$ with conjugate symmetry. Summary~\ref*{subsec:poly-vector-transformation-complex} (\autoref{subsec:poly-vector-transformation-complex}) gives the following canonical embedding:

$\sigma_c: f(X) \in \mathbb{R}[X]/(X^n + 1) \longrightarrow (f(\omega),f(\omega^3),f(\omega^5), \cdots, f(\omega^{2n-1})) \in \mathbb{\hat{C}}^{n} \text{ } (\longrightarrow \mathbb{C}^{\frac{n}{2}})$


, where $\omega = e^{i\pi/n}$ is a root (i.e., the primitive $(\mu=2n)$-th root) of the $(\mu=2n)$-th cyclotomic polynomial $X^n + 1$ defined over complex numbers, and $\mathbb{\hat{C}}^{n}$ is $n$-dimensional complex special vector space whose second-half elements are reverse-ordered conjugates of the first-half elements. And $\mathbb{\hat{C}}^{n}$ is isomorphic to $\mathbb{{C}}^{\frac{n}{2}}$, because the second-half elements of $\mathbb{\hat{C}}^{n}$ are automatically determined by its first-half elements. Therefore, the $\sigma_c$ mapping is essentially an isomorphism between $\dfrac{n}{2}$-slot complex vectors $\vec{v} \in \mathbb{C}^{\frac{n}{2}}$ and $(n-1)$-degree (or lesser degree) real-number polynomials $\mathbb{R}[X] / (X^n + 1)$. Therefore, CKKS' batch encoding scheme encodes an $\dfrac{n}{2}$-slot complex input vector into an $(n-1)$-degree (or lesser degree) real-number polynomial, and the decoding process is a reverse of this.

In addition, remember that in BFV, we updated $W$ and $W^T$ to $\hathat W$ and $\hathat W^*$ (Summary~\ref*{subsubsec:bfv-rotation-summary} in \autoref{subsubsec:bfv-rotation-summary}) to support homomorphic rotation of input vector slots. Likewise, the CKKS batch encoding scheme uses $\hathat W$ and $\hathat W^*$ instead of $W$ and $W^T$ in order to support homomorphic rotation. Therefore, the CKKS batch encoding scheme's isomorphic mapping is updated as follows:


$\sigma_c: f(X) \in \mathbb{R}[X]/(X^n + 1) \longrightarrow \bm ( f(\omega^{J(0)}),f(\omega^{J(1)}),f(\omega^{J(2)}), \cdots, f(\omega^{J(\frac{n}{2} - 1)}), $

\textcolor{white}{$ f(X) \in \mathbb{R}[X]/(X^n + 1) \longrightarrow \bm ( $} $\cdots, f(\omega^{J_*(0)}),f(\omega^{J_*(1)}),f(\omega^{J_*(2)}), \cdots, f(\omega^{J_*(\frac{n}{2} - 1)}) \bm) \in \mathbb{\hat{C}}^{n} \text{ } (\longrightarrow \mathbb{C}^{\frac{n}{2}})$

, where $J(h) = 5^h \bmod 2n$, a rotation helper formula.


$ $

The encoding schemes of BFV and CKKS have the following differences:

\begin{itemize}
\item \textbf{Type of Input Slot Values:} BFV's input slot values are ${n}$ integers modulo $t$, which are encoded into $n$ polynomial coefficients (i.e., modulo-$t$ integers). On the other hand, CKKS's input slot values are $\dfrac{n}{2}$ complex numbers, which are encoded into $n$ polynomial coefficients (i.e., real numbers).

\item \textbf{Type of Polynomial Coefficients:} BFV encodes into $\mathbb{Z}_t[X]/(X^n+1)$. CKKS first interpolates real coefficients and then rounds the scaled coefficients to integers. The encrypted polynomial has integer coefficients in both schemes.

\item \textbf{Scaling Factor:} BFV uses $\Delta=\lfloor q/t\rfloor$ at ciphertext modulus $q$. CKKS chooses $\Delta$ according to the desired precision, message bounds, and modulus chain. Correct interpretation requires a centered no-wrap bound on the whole decryption phase, including all accumulated errors; zero and negative coefficients are allowed. At every level, require $\|\Delta M+E\|_\infty<q_l/2$.

\item \textbf{Encoding Precision:} In the case of BFV, during its decoding process, BFV's down-scaled polynomial coefficients $\dfrac{ \Delta\cdot \vec{m} }{\Delta}$ preserve the precision of input values
%as far as $\Delta$ is sufficiently large enough to round off the scaling error (as explained in \autoref{subsubsec:bfv-enc-dec-decoding1})
. On the other hand, CKKS's down-scaled polynomial coefficients may lose their precision if their original input values have too many decimal digits so that the scaling factor cannot left-shift all of them to make them part of the integer domain, which means that some lower decimal digits of the input value may be rounded off, which loses precision of the original input. For example, suppose the polynomial coefficient $m_i = \dfrac{1}{3} = 0.33333\cdots$, and the scaling factor $\Delta = 100$. Then, the scaled coefficient $\lceil \Delta m_i \rfloor = 33$, and down-scaling it gives $\dfrac{33}{100} = 0.33$. Since $0.33 \neq 0.33333\cdots$, CKKS's encoding and decoding process does not always guarantee exact precision. Due to this encoding error, CKKS is called an \textit{approximate} encryption scheme. The impact of this encoding error can grow over homomorphic operations which increases the magnitude of error and the decoded result would gradually become more deviated from the expected exact value. One way to reduce CKKS's encoding error is to increase $\Delta$, and thereby left-shift more decimal digits to make them part of the scaled integer digits.

\end{itemize}

$ $


\para{Structure of \boldmath$\vec{v}_{'} \in \mathbb{\hat C}^n$:} Note that the original decoding scheme for $\vec{v}_{'}$ described in Summary~\ref*{subsec:poly-vector-transformation-complex} (\autoref{subsec:poly-vector-transformation-complex}) was:

$\vec{v}_{'} = \bm{(} \text{ } M(\omega), \text{ } M(\omega^3), \text{ } M(\omega^5), \cdots, M(\omega^{2n-3}), \text{ } M(\omega^{2n-1}) \bm{)}$

, which decodes to a Hermitian vector:

$\vec{v}_{'} = (v_0, v_1, \cdots, v_{\frac{n}{2} - 1}, \overline v_{\frac{n}{2} - 1}, \cdots, \overline v_1, \overline v_0 )$

, whose second-half elements are reverse-ordered conjugates of the first-half elements.

$ $

However, by replacing $W$ and $W^T$ with $\hathat W$ and $\hathat W^*$, we changed the above decoding scheme to the following that supports homomorphic rotation:


{\small $\vec{v}_{'} =  \bm{(} \text{ }
M(\omega^{J(0)}), \text{ } M(\omega^{J(1)}), \text{ } M(\omega^{J(2)}), \cdots,  M(\omega^{J(\frac{n}{2}-1)}), \text{ } M(\omega^{J_*(0)}), \text{ } M(\omega^{J_*(1)}), \cdots,  M(\omega^{J_*(\frac{n}{2}-1)}) \text{ } \bm{)}$}

$\textcolor{white}{\vec{v}_{'}} =  \bm{(} \text{ }
M(\omega^{J(0)}), \text{ } M(\omega^{J(1)}), \text{ } M(\omega^{J(2)}), \cdots,  M(\omega^{J(\frac{n}{2}-1)}), \text{ } M(\overline\omega^{J(0)}), \text{ } M(\overline\omega^{J(1)}), \cdots,  M(\overline\omega^{J(\frac{n}{2}-1)}) \text{ } \bm{)}$

\textcolor{red}{ $\rhd$ because $\omega^{-1} = (e^{\frac{i\pi}{n}})^{-1} = e^{\frac{-i\pi}{n}} = \overline\omega$, given  $J(h) = 5^h \bmod 2n$, and $J_*(h) = -5^h \bmod 2n$}

, which decodes to a \textit{forward-ordered} (not reverse-ordered) Hermitian vector as follows:

$\vec{v}_{'} = (v_0, v_1, \cdots, v_{\frac{n}{2} - 1}, \overline v_0, \overline v_1, \cdots, \overline v_{\frac{n}{2} - 1})$

, whose second-half elements are conjugates of the first-half elements with the same order. Upon homomorphic rotation (which will be explained in \autoref{subsec:ckks-rotation}), just like in BFV's homomorphic rotation, the first-half elements and the second-half elements of $\vec{v}_{'}$ rotate within their own group in a wrapping manner.



We summarize CKKS's encoding and decoding procedure as follows, which is similar to BFV's encoding and decoding procedure (described in Summary~\ref*{subsubsec:bfv-encoding-summary} in \autoref{subsubsec:bfv-encoding-summary}):

\begin{tcolorbox}[title={\textbf{\tboxlabel{\ref*{subsec:ckks-encoding-decoding}} CKKS's Encoding and Decoding}}]

\textbf{\underline{Input}:} An $\dfrac{n}{2}$-slot complex vector $\vec{v} = (v_0, v_1, \cdots, v_{\frac{n}{2}-1}) \in \mathbb{C}^{\frac{n}{2}}$

\par\noindent\rule{\textwidth}{0.4pt}

\textbf{\underline{Encoding}:}

$ $

\begin{enumerate}
\item Convert (i.e., isomorphically transform) $\vec{v}$ into an $n$-slot \textit{forward-ordered} Hermitian vector $\vec{v}_{'}$ as follows:

$\vec{v}_{'} = (v_0, v_1, \cdots, v_{\frac{n}{2}-1}, \overline v_0, \overline v_1, \cdots, \overline v_{\frac{n}{2}-1}) \in \mathbb{\hat C}^{n}$

\item Convert $\vec{v}_{'}$ into a real number vector $\vec{m}$ by applying the transformation $\vec{m} = \dfrac{\hathat W \cdot I_n^R \cdot \vec{v}_{'}}{n}$

, where $\hathat{W}$ is a basis of the $n$-dimensional vector space crafted as follows:

$ $

\noindent{\footnotesize{\noindent $\noindent\hathat W = \begin{bmatrix}
1 & 1 & \cdots & 1 & 1 & 1 & \cdots & 1\\
(\omega^{J(\frac{n}{2} - 1)}) & (\omega^{J(\frac{n}{2} - 2)}) & \cdots & (\omega^{J(0)}) & (\omega^{J_*(\frac{n}{2} - 1)}) & (\omega^{J_*(\frac{n}{2} - 2)}) & \cdots & (\omega^{J_*(0)})\\
(\omega^{J(\frac{n}{2} - 1)})^2 & (\omega^{J(\frac{n}{2} - 2)})^2 & \cdots & (\omega^{J(0)})^2 & (\omega^{J_*(\frac{n}{2} - 1)})^2 & (\omega^{J_*(\frac{n}{2} - 2)})^2 & \cdots & (\omega^{J_*(0)})^2 \\
\vdots & \vdots & \ddots & \vdots & \vdots & \ddots & \vdots & \vdots \\
(\omega^{J(\frac{n}{2} - 1)})^{n-1} & (\omega^{J(\frac{n}{2} - 2)})^{n-1} & \cdots & (\omega^{J(0)})^{n-1} & (\omega^{J_*(\frac{n}{2} - 1)})^{n-1} & (\omega^{J_*(\frac{n}{2} - 2)})^{n-1} & \vdots  & (\omega^{J_*(0)})^{n-1}
\end{bmatrix}$}}

\textcolor{red}{ $\rhd$ where $\omega = e^{i\pi/n} = \cos \left(\dfrac{\pi}{n}\right) + i\sin \left(\dfrac{\pi}{n}\right)$,  $J(h) = 5^h \bmod 2n$,  and  $J_*(h) = -5^h \bmod 2n$}

\small{\noindent $ = \begin{bmatrix}
1 & 1 & \cdots & 1 & 1 & 1 & \cdots & 1\\
(\omega^{J(\frac{n}{2} - 1)}) & (\omega^{J(\frac{n}{2} - 2)}) & \cdots & (\omega^{J(0)}) & (\overline\omega^{J(\frac{n}{2} - 1)}) & (\overline\omega^{J(\frac{n}{2} - 2)}) & \cdots & (\overline\omega^{J(0)})\\
(\omega^{J(\frac{n}{2} - 1)})^2 & (\omega^{J(\frac{n}{2} - 2)})^2 & \cdots & (\omega^{J(0)})^2 & (\overline\omega^{J(\frac{n}{2} - 1)})^2 & (\overline\omega^{J(\frac{n}{2} - 2)})^2 & \cdots & (\overline\omega^{J(0)})^2 \\
\vdots & \vdots & \ddots & \vdots & \vdots & \ddots & \vdots & \vdots \\
(\omega^{J(\frac{n}{2} - 1)})^{n-1} & (\omega^{J(\frac{n}{2} - 2)})^{n-1} & \cdots & (\omega^{J(0)})^{n-1} & (\overline\omega^{J(\frac{n}{2} - 1)})^{n-1} & (\overline\omega^{J(\frac{n}{2} - 2)})^{n-1} & \vdots  & (\overline\omega^{J(0)})^{n-1}
\end{bmatrix}$}

\textcolor{red}{ $\rhd$ because $\omega^{-1} = e^{\frac{-i\pi}{n}} = \overline{e^{\frac{i\pi}{n}}} = \overline\omega$}

$ $

\item Convert $\vec{m}$ into a scaled integer vector $\lceil \Delta\vec{m} \rfloor \approx \Delta \vec{m}$, where $\Delta$ is a scaling factor bigger than 1 such that $\Delta m_i$ never overflows or underflows $q_0$ (i.e., $0 \leq \Delta m_i < q_0$ or $-\dfrac{q_0}{2} \leq \Delta m_i < \dfrac{q_0}{2}$) in all cases, even across all homomorphic operations. The finally encoded plaintext polynomial is $\Delta M = \sum\limits_{i=0}^{n-1} \lceil \Delta m_i \rfloor X^i \text{ } \in \mathbb{Z}_q[X] / (X^n + 1)$. The rounding process of $\lceil \Delta \vec{m} \rfloor$ during the encoding process causes an encoding error, which makes CKKS an approximate encryption scheme.

\end{enumerate}

\par\noindent\rule{\textwidth}{0.4pt}

\textbf{\underline{Decoding}: } From the plaintext polynomial $\Delta M = \sum\limits_{i=0}^{n-1}\Delta m_iX^i$, recover $\vec{m} = \dfrac{\Delta \vec{m}}{\Delta}$. Then,
compute $\vec{v}_{'} = \hathat W^* \cdot \vec{m}$, where:

\noindent $\hathat{W}^* = \begin{bmatrix}
1 & (\omega^{J(0)}) & (\omega^{J(0)})^2 & \cdots & (\omega^{J(0)})^{n-1}\\
1 & (\omega^{J(1)}) & (\omega^{J(1)})^2 & \cdots & (\omega^{J(1)})^{n-1}\\
1 & (\omega^{J(2)}) & (\omega^{J(2)})^2 & \cdots & (\omega^{J(2)})^{n-1}\\
\vdots & \vdots & \vdots & \ddots & \vdots \\
1 & (\omega^{J(\frac{n}{2}-1)}) & (\omega^{J(\frac{n}{2}-1)})^2 & \cdots & (\omega^{J(\frac{n}{2}-1)})^{n-1}\\
1 & (\omega^{J_*(0)}) & (\omega^{J_*(0)})^2 & \cdots & (\omega^{J_*(0)})^{n-1}\\
1 & (\omega^{J_*(1)}) & (\omega^{J_*(1)})^2 & \cdots & (\omega^{J_*(1)})^{n-1}\\
1 & (\omega^{J_*(2)}) & (\omega^{J_*(2)})^2 & \cdots & (\omega^{J_*(2)})^{n-1}\\
\vdots & \vdots & \vdots & \ddots & \vdots \\
1 & (\omega^{J_*(\frac{n}{2}-1)}) & (\omega^{J_*(\frac{n}{2}-1)})^2 & \cdots & (\omega^{J_*(\frac{n}{2}-1)})^{n-1}\\
\end{bmatrix}$

$ $

\noindent $ = \begin{bmatrix}
1 & (\omega^{J(0)}) & (\omega^{J(0)})^2 & \cdots & (\omega^{J(0)})^{n-1}\\
1 & (\omega^{J(1)}) & (\omega^{J(1)})^2 & \cdots & (\omega^{J(1)})^{n-1}\\
1 & (\omega^{J(2)}) & (\omega^{J(2)})^2 & \cdots & (\omega^{J(2)})^{n-1}\\
\vdots & \vdots & \vdots & \ddots & \vdots \\
1 & (\omega^{J(\frac{n}{2}-1)}) & (\omega^{J(\frac{n}{2}-1)})^2 & \cdots & (\omega^{J(\frac{n}{2}-1)})^{n-1}\\
1 & (\overline\omega^{J(0)}) & (\overline\omega^{J(0)})^2 & \cdots & (\overline\omega^{J(0)})^{n-1}\\
1 & (\overline\omega^{J(1)}) & (\overline\omega^{J(1)})^2 & \cdots & (\overline\omega^{J(1)})^{n-1}\\
1 & (\overline\omega^{J(2)}) & (\overline\omega^{J(2)})^2 & \cdots & (\overline\omega^{J(2)})^{n-1}\\
\vdots & \vdots & \vdots & \ddots & \vdots \\
1 & (\overline\omega^{J(\frac{n}{2}-1)}) & (\overline\omega^{J(\frac{n}{2}-1)})^2 & \cdots & (\overline\omega^{J(\frac{n}{2}-1)})^{n-1}\\
\end{bmatrix}$

$ $

$ $

, and extract only the first $\dfrac{n}{2}$ elements in the \textit{forward-ordered} Hermitian vector $\vec{v}_{'}$ to recover the input vector $\vec{v}$.

\end{tcolorbox}

\begin{comment}
\begin{table}[h] %usepackage{array}
\centering
\begin{tabular}{|>{\centering\arraybackslash}p{0.2\columnwidth}||>{\centering\arraybackslash}p{0.75\columnwidth}|}
\hline \hline
& \textbf{Encoded Polynomial $M(X) \in \mathcal{R}_{\langle n \rangle}$ over $X \in \mathbb{C}$ (complex numbers)} \\ \hline \hline
\textbf{Range of Input Slot Value}& Depends on the range of $M(X)$'s coefficient\\ \hline
\textbf{Range of $\bm{M(X)}$'s coefficients}&Any real number $m_i$ with the custom precision scaling factor $\Delta$ under the constraint: $0 \leq \Delta m_i \leq q-1$ (or $-\dfrac{q}{2} \leq \Delta m_i \leq \dfrac{q}{2} - 1$) \\\hline
%\textbf{Range of $\bm Y$}& $\left[0, \text{ } (1+i)\cdot {t}\cdot 2^{ \lceil \frac{n+1}{2} \rceil}\right]$\\
%& or $\left[-(1+i)\cdot \dfrac{t}{2}\cdot 2^{ \lceil \frac{n + 1}{2} \rceil}, \text{ } (1+i)\cdot \dfrac{t}{2}\cdot 2^{ \lceil \frac{n+1}{2} \rceil}\right]$ (a rough estimate)\\
\hline
\end{tabular}
\caption{The encoding capacity of $M(X) \in \mathcal{R}_{\langle n, t \rangle}$.}
\label{tab:polynomial-encoding-capacity-complex}
\end{table}

\para{CKKS's Encoding Range:} \autoref{tab:polynomial-encoding-capacity-complex} depicts the range of input values that can be encoded by the polynomials $M(X) \in \mathcal{R}_{\langle n \rangle}$ over $X \in \mathbb{C}$. As illustrated in \autoref{tab:polynomial-encoding-capacity-complex}, the encoding range of complex number inputs is determined by the size of the $n$ and $q$ parameters. Therefore, these two parameters should be chosen sufficiently large to encode all values that can be used by an application.
\end{comment}

\para{CKKS's Approximation Property:} Let $P=\lceil\Delta M\rfloor$ be coefficient-wise nearest-integer rounding. Then $P=\Delta M+E_{\mathrm{enc}}$ with $\|E_{\mathrm{enc}}\|_\infty\leq 1/2$. Each decoded slot has encoding error at most $n/(2\Delta)$ by the triangle inequality at unit-modulus roots. This conservative coefficient-norm bound is distinct from the canonical embedding norm used in CKKS analyses. Encryption and later computation add further errors; approximate decoding does not automatically remove them.


\subsubsection{Example}
\label{subsubsec:ckks-encoding-ex}


Suppose our input complex vector's dimension $\dfrac{n}{2} = 2$, the bounding polynomial degree $n$ = 4, and the scaling factor $\Delta = 1024$.

Our basis of the $n$-dimensional vector space

$\hathat{W}= \begin{bmatrix}
1 & 1 & 1 & 1\\
\omega^{J(1)} & \omega^{J(0)} & \overline{\omega}^{J(1)} & \overline{\omega}^{J(0)}\\
(\omega^{J(1)})^2 & ({\omega^{J(0)}})^2 & (\overline{\omega}^{J(1)})^2 & (\overline{\omega}^{J(0)})^2\\
(\omega^{J(1)})^3 & (\omega^{J(0)})^3 & (\overline{\omega}^{J(1)})^3 & (\overline{\omega}^{J(0)})^3\\
\end{bmatrix}$
$= \begin{bmatrix}
1 & 1 & 1 & 1\\
\omega^5 & \omega & \overline{\omega}^5 & \overline{\omega}\\
\omega^2 & \omega^2 & \overline{\omega^2} & \overline{\omega}^2\\
\omega^7 & \omega^3 & \overline{\omega}^7 & \overline{\omega}^3\\
\end{bmatrix}$

, where $\omega = e^{i\pi/n} = \cos \left(\dfrac{\pi}{n}\right) + i\sin \left(\dfrac{\pi}{n}\right)$

$ $


Given this setup, suppose we have the input complex vector $\vec{v} = (1.1 + 4.3i, \text{ } 3.5 - 1.4i)$ to encode.

$ $

First, construct the forward-ordered Hermitian vector $\vec{v}_{'} = (1.1 + 4.3i, \text{ } 3.5 - 1.4i, \text{ } 1.1 - 4.3i, \text{ } 3.5 + 1.4i)$.

$ $

Next, convert the complex vector $\vec{v}_{'}$ into a real number vector $\vec{m}$ by applying the transformation:


$\vec{m} = \dfrac{\hathat{W} \cdot I_n^R \cdot \vec{v}_{'}}{n} = \dfrac{1}{4} \cdot \begin{bmatrix}
1 & 1 & 1 & 1\\
\omega^5 & \omega & \overline{\omega}^5 & \overline{\omega}\\
\omega^2 & \omega^2 & \overline{\omega}^2 & \overline{\omega}^2\\
\omega^7 & \omega^3 & \overline{\omega}^7 & \overline{\omega}^3\\
\end{bmatrix} \cdot
\begin{bmatrix}
0 & 0 & 0 & 1 \\
0 & 0 & 1 & 0 \\
0 & 1 & 0 & 0 \\
1 & 0 & 0 & 0
\end{bmatrix}
\cdot
\begin{bmatrix} 1.1 + 4.3i\\3.5 - 1.4i\\1.1 - 4.3i\\3.5 + 1.4i \end{bmatrix}$

$ $

$= \dfrac{\hathat W \cdot I_n^R \cdot \vec{v}_{'}}{n} = \dfrac{1}{4} \cdot\begin{bmatrix}
1 & 1 & 1 & 1\\
\overline{\omega} & \overline{\omega}^5 & \omega  & \omega^5\\
\overline{\omega}^2 & \overline{\omega}^2 & \omega^2 & \omega^2\\
\overline{\omega}^3 & \overline{\omega}^7 & \omega^3 & \omega^7\\
\end{bmatrix}
\cdot
\begin{bmatrix} 1.1 + 4.3i\\3.5 - 1.4i\\1.1 - 4.3i\\3.5 + 1.4i \end{bmatrix}$

$ $

$= \dfrac{1}{4} \cdot
\begin{bmatrix}
(1.1 + 4.3i) + (3.5 - 1.4i) + (1.1 - 4.3i) + (3.5 + 1.4i)\\
(1.1 + 4.3i)\overline{\omega} + (3.5 - 1.4i)\overline{\omega}^5 + (1.1 - 4.3i)\omega+ (3.5 + 1.4i)\omega^5 \\
(1.1 + 4.3i)\overline{\omega^2} + (3.5 - 1.4i)\overline{\omega^2} + (1.1 - 4.3i)\omega^2+ (3.5 + 1.4i)\omega^2 \\
(1.1 + 4.3i)\overline{\omega}^3 + (3.5 - 1.4i)\overline{\omega}^7 + (1.1 - 4.3i)\omega^3 + (3.5 + 1.4i)\omega^7 \\
\end{bmatrix}$

$ $

$= \dfrac{1}{4} \cdot
\begin{bmatrix}
(1.1 + 4.3i) + (3.5 - 1.4i) + (1.1 - 4.3i) + (3.5 + 1.4i)\\
(1.1 + 4.3i)\overline{\omega} + (3.5 - 1.4i)\overline{\omega}^5 + (1.1 - 4.3i)\omega+ (3.5 + 1.4i)\omega^5 \\
(1.1 + 4.3i)\overline{\omega^2} + (3.5 - 1.4i)\overline{\omega^2} + (1.1 - 4.3i)\omega^2+ (3.5 + 1.4i)\omega^2 \\
(1.1 + 4.3i)\overline{\omega}^3 + (3.5 - 1.4i)\overline{\omega}^7 + (1.1 - 4.3i)\omega^3 + (3.5 + 1.4i)\omega^7 \\
\end{bmatrix}$

$ $

$= \dfrac{1}{4} \cdot
\begin{bmatrix}
9.2\\
1.1(\overline{\omega} + \omega) + 4.3i(\overline{\omega} - \omega) +
3.5(\overline{\omega}^5 + \omega^5) - 1.4i(\overline{\omega}^5 - \omega^5)\\
1.1(\overline{\omega}^2 + {\omega}^2) + 4.3i(\overline{\omega}^2 - {\omega}^2) +
3.5(\overline{\omega}^2 + {\omega}^2) - 1.4i(\overline{\omega}^2 - {\omega}^2)\\
1.1(\overline{\omega}^3 + {\omega}^3) + 4.3i(\overline{\omega}^3 - {\omega}^3) +
3.5(\overline{\omega}^7 + {\omega}^7) - 1.4i(\overline{\omega}^7 - {\omega}^7)\\
\end{bmatrix}$

$ $

$= \dfrac{1}{4} \cdot
\begin{bmatrix}
9.2\\
1.1\left(2\cos\dfrac{\pi}{4}\right) - 4.3i\left(2i\sin\dfrac{\pi}{4}\right) + 3.5\left(2\cos\dfrac{5\pi}{4}\right) + 1.4i\left(2i\sin\dfrac{5\pi}{4}\right)\\
1.1\left(2\cos\dfrac{\pi}{2}\right) - 4.3i\left(2i\sin\dfrac{\pi}{2}\right) + 3.5\left(2\cos\dfrac{\pi}{2}\right) + 1.4i\left(2i\sin\dfrac{\pi}{2}\right)\\
1.1\left(2\cos\dfrac{3\pi}{4}\right) - 4.3i\left(2i\sin\dfrac{3\pi}{4}\right) + 3.5 \left(2\cos\dfrac{7\pi}{4}\right) + 1.4i\left(2i\sin\dfrac{7\pi}{4}\right)\\
\end{bmatrix}$

$ $

$= \dfrac{1}{4} \cdot
\begin{bmatrix}
9.2\\
1.1\left(2\dfrac{\sqrt{2}}{2}\right) + 4.3\left(2\dfrac{\sqrt{2}}{2}\right) + 3.5\left(-2\dfrac{\sqrt{2}}{2}\right) - 1.4\left(-2\dfrac{\sqrt{2}}{2}\right)\\
1.1(2\cdot 0) + 4.3(2\cdot1) + 3.5(2\cdot 0) - 1.4(2\cdot 1)\\
1.1\left(-2\dfrac{\sqrt{2}}{2}\right) + 4.3\left(2\dfrac{\sqrt{2}}{2}\right) + 3.5 \left(2\dfrac{\sqrt{2}}{2}\right) - 1.4\left(-2\dfrac{\sqrt{2}}{2}\right)\\
\end{bmatrix}$

$ $

$= 0.25 \cdot
\begin{bmatrix}
9.2\\
1.1\sqrt{2} + 4.3{\sqrt{2}} - 3.5\sqrt{2} + 1.4\sqrt{2}\\
1.1(0) + 4.3(2) - 3.5(0) - 1.4(2)\\
-1.1\sqrt{2} + 4.3\sqrt{2} + 3.5\sqrt{2} + 1.4\sqrt{2}\\
\end{bmatrix} =
\begin{bmatrix}
2.3\\
0.825\sqrt{2}\\
1.45\\
2.025\sqrt{2}\\
\end{bmatrix} \approx (2.3, \text{ } 1.1667, \text{ } 1.45, \text{ } 2.8638)$

$ $

Convert the real number vector $\vec{m}$ into a scaled integer vector $\Delta\vec{m}$ by $\Delta$-scaling and rounding as follows:

$\Delta\vec{m} \approx \lceil \Delta \vec{m} \rfloor = \lceil 1024 \cdot (2.3, \text{ } 1.1667, \text{ } 1.45, \text{ } 2.8638) \rfloor = (2355, \text{ } 1195, \text{ } 1485, \text{ } 2933)$

$ $

Finally, $\vec{v} = (1.1 + 4.3i, \text{ } 3.5 - 1.4i)$ has been encoded into the plaintext polynomial $M(X)$ as follows:

$\Delta M(X) = 2355 +  1195X + 1485X^2 + 2933X^3 \in \mathcal{R}_{\langle 4 \rangle}$

$ $


To decode $\vec{m}$, we compute:

$\vec{v}_{'} = \dfrac{\hathat W^*\cdot\Delta\vec{m}}{\Delta} = \begin{bmatrix}
1,\omega,\omega^2, \omega^3\\
1,\omega^5, \omega^2,\omega^7\\
1,\overline{\omega}, \overline{\omega^2}, \overline{\omega^3}\\
1,\overline{\omega^5}, \overline{\omega^2}, \overline{\omega^7}
\end{bmatrix}$
$\cdot \begin{bmatrix}
2355\\1195\\1485\\2933
\end{bmatrix}\cdot \dfrac{1}{1024}$

$=
\begin{bmatrix}
1,\dfrac{\sqrt{2}}{2} + \dfrac{i\sqrt{2}}{2},i, -\dfrac{\sqrt{2}}{2} + \dfrac{i\sqrt{2}}{2}\\
1,-\dfrac{\sqrt{2}}{2} - \dfrac{i\sqrt{2}}{2}, i,\dfrac{\sqrt{2}}{2} - \dfrac{i\sqrt{2}}{2}\\
1,\dfrac{\sqrt{2}}{2} - \dfrac{i\sqrt{2}}{2}, -i, -\dfrac{\sqrt{2}}{2} - \dfrac{i\sqrt{2}}{2}\\
1,-\dfrac{\sqrt{2}}{2} + \dfrac{i\sqrt{2}}{2}, -i, \dfrac{\sqrt{2}}{2} + \dfrac{i\sqrt{2}}{2}
\end{bmatrix}$
$\cdot \begin{bmatrix}
2.2998046875\\1.1669921875\\1.4501953125\\2.8642578125
\end{bmatrix}$

$ $

$ \approx (1.0997+4.3007i, \text{ } 3.5000-1.4003i, \text{ } 1.0997-4.3007i, \text{ } 3.5000+1.4003i)$

$ $

Extract the first $\dfrac{n}{2} = 2$ elements in the Hermitian vector $\vec{v}_{'}$ to recover the input vector:

$(1.0997+4.3007i, \text{ } 3.5000-1.4003i)$

$\approx (1.1 + 4.3i, \text{ } 3.5 - 1.4i) = \vec{v}$ \textcolor{red}{  $\rhd$ The original input vector}

$ $

The integer rounding of the scaled inverse embedding produces the small difference between the original and decoded slots. The displayed decimal intermediate coefficients are approximations; the integer encoding is computed from the exact expressions involving $\sqrt2$.

$ $

\para{Source Code:} Examples of CKKS encoding can be executed by running \href{https://github.com/wang2000star/fhe-textbook-cn/blob/chinese-edition/source%20code/ckks.py}{\underline{this Python script}}.



\subsection{Encryption and Decryption}
\label{subsec:ckks-enc-dec}

CKKS encryption operates on $P=\lceil\Delta M\rfloor\in R$, not on arbitrary real coefficients directly. The following is the symmetric-key RLWE form; public-key CKKS encryption uses a public RLWE sample and has a corresponding combined error.

\begin{tcolorbox}[title={\textbf{\tboxlabel{\ref*{subsec:ckks-enc-dec}} CKKS Encryption and Decryption}}]
Choose a small secret $S\in R$, sample $A\leftarrow R_q$ uniformly and a small integer error $E_0$, and output
\[
B=[-AS+P+E_0]_q,\qquad \textsf{ct}=(A,B)\in R_q^2.
\]
Define the raw decryption phase by
\[
\mathsf{Phase}_S(\textsf{ct})=[B+AS]_q=P+E_0=\Delta M+E,
\qquad E=E_{\mathrm{enc}}+E_0,
\]
provided $\|\Delta M+E\|_\infty<q/2$. Here $[\cdot]_q$ selects centered coefficient representatives. Decode by evaluating the phase at the selected roots and dividing by $\Delta$:
\[
\widehat v_j=\Delta^{-1}\mathsf{Phase}_S(\textsf{ct})(\omega^{J(j)})
=v_j+\Delta^{-1}E(\omega^{J(j)}).
\]
Thus $|\widehat v_j-v_j|\leq\|E\|_{\mathrm{can},\infty}/\Delta\leq n\|E\|_\infty/\Delta$. The condition $|e_i|<\Delta/2$ alone does not imply slot precision $1/\Delta$. Rounding to multiples of $1/\Delta$ cannot remove general CKKS noise.
\end{tcolorbox}

Below, $\textsf{RLWE}^{-1}$ denotes this raw phase when used in algebraic derivations; division by the tracked scale and slot decoding are separate operations. $\textsf{RLWE}_{S,\sigma}(P)$ denotes an encryption relation, not a claim that derived ciphertexts retain the original independent error distribution $\chi_\sigma$.

\subsection{Ciphertext-to-Ciphertext Addition}
\label{subsec:ckks-add-cipher}

CKKS adds ciphertext components as in BFV (Summary~\ref*{subsec:bfv-add-cipher}, \autoref{subsec:bfv-add-cipher}). The ciphertexts must have compatible keys, moduli, and scales. The resulting phase error is $E^{\langle1\rangle}+E^{\langle2\rangle}$, and the summed phase must satisfy the no-wrap bound.




\begin{tcolorbox}[title={\textbf{\tboxlabel{\ref*{subsec:ckks-add-cipher}} CKKS Ciphertext-to-Ciphertext Addition}}]
$\textsf{RLWE}_{S, \sigma}(\Delta M^{\langle 1 \rangle} ) + \textsf{RLWE}_{S, \sigma}(\Delta M^{\langle 2 \rangle} ) $

$ = ( A^{\langle 1 \rangle}, \text{ } B^{\langle 1 \rangle}) + (A^{\langle 2 \rangle}, \text{ } B^{\langle 2 \rangle}) $

$ = ( A^{\langle 1 \rangle} + A^{\langle 2 \rangle}, \text{ } B^{\langle 1 \rangle} + B^{\langle 2 \rangle} ) $

$= \textsf{RLWE}_{S, \sigma}(\Delta(M^{\langle 1 \rangle} + M^{\langle 2 \rangle}) )$
\end{tcolorbox}


\subsection{Ciphertext-to-Plaintext Addition}
\label{subsec:ckks-add-plain}

CKKS adds a plaintext to the $B$ component as in BFV (Summary~\ref*{subsec:bfv-add-plain}, \autoref{subsec:bfv-add-plain}). Encode the plaintext at the ciphertext's actual scale $\Delta$, and include its encoding error in the error budget. In the following shorthand, $\Delta\Lambda$ denotes the rounded integer encoding.



\begin{tcolorbox}[title={\textbf{\tboxlabel{\ref*{subsec:ckks-add-plain}} CKKS Ciphertext-to-Plaintext Addition}}]
$\textsf{RLWE}_{S, \sigma}(\Delta M) + \Delta\Lambda $

$=  (A, \text{ } B) + \Delta\Lambda$

$=  (A, \text{ } B + \Delta\cdot\Lambda)$

$= \textsf{RLWE}_{S, \sigma}(\Delta (M + \Lambda) )$
\end{tcolorbox}




\subsection{Ciphertext-to-Ciphertext Multiplication}
\label{subsec:ckks-mult-cipher}

CKKS multiplication first tensors ciphertexts at their current modulus, then usually relinearizes and rescales. Tensoring changes the scale from $\Delta_1,\Delta_2$ to $\Delta_1\Delta_2$; rescaling by $w_l=q_l/q_{l-1}$ changes it to $\Delta_1\Delta_2/w_l$ and consumes one level. The following derivation uses the sign convention $B=-AS+\Delta M+E\pmod q$, as in Theorem~\ref*{subsec:glwe-alternative} (\autoref{subsec:glwe-alternative}). We first derive the exact noisy phase, then explain two relinearization methods and rescaling.


\subsubsection{Synthetic Ciphertext Derivation}
\label{subsubsec:ckks-mult-cipher-relation}

For two ciphertexts at the same modulus and secret key, let their phases be $\Delta M^{\langle 1\rangle}+E^{\langle 1\rangle}$ and $\Delta M^{\langle 2\rangle}+E^{\langle 2\rangle}$. Define, modulo $q$,
\[
D_0=B^{\langle 1\rangle}B^{\langle 2\rangle},\quad
D_1=A^{\langle 1\rangle}B^{\langle 2\rangle}+A^{\langle 2\rangle}B^{\langle 1\rangle},\quad
D_2=A^{\langle 1\rangle}A^{\langle 2\rangle}.
\]
The exact phase identity is
\[
D_0+D_1S+D_2S^2\equiv\Delta^2M^{\langle 1\rangle}M^{\langle 2\rangle}+F\pmod q,
\]
\[
F=\Delta M^{\langle 1\rangle}E^{\langle 2\rangle}
+\Delta M^{\langle 2\rangle}E^{\langle 1\rangle}
+E^{\langle 1\rangle}E^{\langle 2\rangle}.
\]
CKKS decoding retains these error terms. Relinearization forms $C_\alpha=(D_1,D_0)$ and a pair $C_\beta$ with phase $D_2S^2+E_{\mathrm{ks}}$. Then $C_\alpha+C_\beta\in R_q^2$ has phase $\Delta^2M^{\langle 1\rangle}M^{\langle 2\rangle}+F+E_{\mathrm{ks}}$. It does not restore the scale to $\Delta$; that is the purpose of the later rescaling step. Without relinearization, ciphertext length grows with the degree in $S$; repeated squaring doubles that degree at every step.

\subsubsection{Relinearization  Method 1 -- Ciphertext Decomposition}
\label{subsubsec:relinearization-gadget-decomposition}

Use an explicit integer gadget for arbitrary $q$: choose an integer $\beta\geq 2$, let $l=\lceil\log_\beta q\rceil$, and set $g_i=\beta^{i-1}$ for $1\leq i\leq l$. Decompose each coefficient of $D_2$ using its unsigned representative into digits $d_i\in\{0,\ldots,\beta-1\}$, so
\[
D_2\equiv\sum_{i=1}^l d_i g_i\pmod q.
\]
This ascending gadget is a separate convention from the descending $q/\beta^i$ weights: the same digit vector cannot be used with both conventions. Generate evaluation-key pairs $K_i=(A_i,B_i)$ satisfying
\[
B_i+A_iS\equiv g_iS^2+E_i\pmod q.
\]
Then
\[
C_\beta=\sum_{i=1}^l d_iK_i,\qquad
\mathsf{Phase}_S(C_\beta)\equiv D_2S^2+\sum_{i=1}^l d_iE_i\pmod q.
\]
For the negacyclic coefficient norm, $\|\sum_i d_iE_i\|_\infty\leq n(\beta-1)\sum_i\|E_i\|_\infty$. Balanced digits or a truncated decomposition require their corresponding digit and residual bounds. The evaluation keys encrypt gadget multiples of $S^2$, without the CKKS message scale $\Delta$.

\subsubsection{Relinearization Method 2 -- Ciphertext Modulus Switch}
\label{subsubsec:relinearization-modulus-switch}

An alternative relinearization method uses an auxiliary positive integer $g$ and an evaluation key at modulus $gq$. Choose $g$ according to the key-switching error budget; it need not be random or a power of two. Sample $A'$ uniformly in $R_{gq}$ and a small integer error $E'$, and set
\[
\mathit{evk}_g=(A',B'),\qquad B'=[-A'S+gS^2+E']_{gq}.
\]
Lift the coefficients of $D_2\in R_q$ to fixed centered integers. Compute the pair $Z=[D_2\mathit{evk}_g]_{gq}$ and set
\[
C_\beta=\left[\left\lceil Z/g\right\rfloor\right]_q.
\]
Write $R_A=\lceil Z_A/g\rfloor-Z_A/g$ and $R_B=\lceil Z_B/g\rfloor-Z_B/g$. Each coefficient has magnitude at most $1/2$. Modulo $q$, the output phase is
\[
\mathsf{Phase}_S(C_\beta)\equiv D_2S^2+D_2E'/g+R_B+R_AS\pmod q.
\]
Hence a coefficient-norm bound on the added error is
\[
\|E_{\mathrm{ks}}\|_\infty\leq
\frac{n\|D_2\|_\infty\|E'\|_\infty}{g}
+\frac{1+\|S\|_1}{2}.
\]
The multiple of $gq$ introduced by reducing $Z$ becomes a multiple of $q$ after division, which disappears modulo $q$. Rounding does not distribute over sums or commute with multiplication by $S$; its contribution is the explicit term $R_B+R_AS$. Large $g$ suppresses the evaluation-key error but does not eliminate this rounding term. The choice of $g$ also enters the security parameters of the extended-modulus evaluation key.


\subsubsection{Rescaling}
\label{subsubsec:ckks-mult-cipher-rescale}

Let $q_l=w_lq_{l-1}$ and let a ciphertext before rescaling have phase $\Delta_{\mathrm{in}}M+F$. Rescaling is coefficient-wise nearest-integer division of centered ciphertext representatives by $w_l$:
\[
\widehat A=[\lceil A/w_l\rfloor]_{q_{l-1}},\qquad
\widehat B=[\lceil B/w_l\rfloor]_{q_{l-1}}.
\]
With $R_A=\lceil A/w_l\rfloor-A/w_l$ and $R_B=\lceil B/w_l\rfloor-B/w_l$,
\[
\widehat B+\widehat AS\equiv
\frac{\Delta_{\mathrm{in}}}{w_l}M+
\frac{F}{w_l}+R_B+R_AS\pmod {q_{l-1}}.
\]
Therefore the actual output scale is $\Delta_{\mathrm{out}}=\Delta_{\mathrm{in}}/w_l$, without rounding the scale metadata. After multiplying ciphertexts of scales $\Delta_1,\Delta_2$, this becomes $\Delta_{\mathrm{out}}=\Delta_1\Delta_2/w_l$. If $\Delta_1=\Delta_2=\Delta$ and $w_l\approx\Delta$, the result is close to $\Delta$, but generally is not equal to it. Relabeling the scale as exactly $\Delta$ introduces a message-dependent error and requires a separate accuracy analysis.

The new phase error is $\widehat F=F/w_l+R_B+R_AS$, not coefficient-wise $\lceil F/w_l\rfloor$. A deterministic bound is
\[
\|\widehat F\|_\infty\leq\|F\|_\infty/w_l+(1+\|S\|_1)/2.
\]
Correct interpretation also requires the complete output phase to lie in the centered interval modulo $q_{l-1}$.

For an RNS chain $q_l=\prod_{i=0}^l w_i$, the primes are chosen for efficient arithmetic and the desired scale schedule. RNS changes the representation, not the total information in $q_l$. RLWE security depends on the dimension, total modulus, secret distribution, error distribution, and evaluation-key parameters; it is not determined only by the smallest and largest RNS primes.

\para{Initial Scaling Factor $\Delta$ Upon Encryption:} Choose $\Delta$ for the desired precision and available modulus headroom, with $q_0$ large enough for the intended final message and error. Setting $\Delta$ near the rescaling primes is convenient, but $\Delta=w_L$ is not mandatory.

\para{Noise Growth:} Multiplication propagates both message-error cross terms and the error product, and rescaling adds a new rounding error. Bounds depend on message and error magnitudes; the error need not increase for every input, for example when multiplying by zero. Rescaling alone does not restore lost numerical precision.

\subsubsection{Comparing BFV and CKKS Bootstrapping}
\label{subsubsec:bfv-bootstrapping-ckks-comparison}

Both schemes refresh a ciphertext by evaluating a suitable decryption-related function, but their exact steps and plaintext semantics differ. In the CKKS procedure below, lifting the ciphertext coefficients exposes the raw phase $\Delta M+E+q_0K$ at a larger modulus; no separate encrypted-secret inner product is needed. Linear coefficient/slot transforms and an approximate modular-reduction function remove the $q_0K$ contribution approximately. BFV bootstrapping instead evaluates an exact finite-plaintext decryption function, with its parameter and digit-extraction conditions.

\para{Multiplicative Level:} BFV at a fixed ciphertext modulus still has a finite noise budget. In the standard CKKS multiply--relinearize--rescale schedule, rescaling consumes a modulus-chain level. Reaching $q_0$ means that no further scheduled rescaling is available; it is not an algebraic prohibition on every multiplication. Bootstrapping can replenish the available levels, while consuming some levels itself (\autoref{subsec:ckks-bootstrapping}).

\para{Limitation of Noise Handling:} Rescaling divides the existing phase error and modulus by the same factor, and adds rounding error. It therefore does not by itself improve the old error-to-modulus ratio or recover lost numerical precision. CKKS bootstrapping raises the available modulus but retains the previous numerical error and adds its own approximation and evaluation errors.


\subsubsection{Summary}
\label{subsubsec:ckks-mult-cipher-summary}

\begin{tcolorbox}[title={\textbf{\tboxlabel{\ref*{subsubsec:ckks-mult-cipher-summary}} CKKS Ciphertext-to-Ciphertext Multiplication}}]
For ciphertexts with the same key and modulus $q_l$, and scales $\Delta_1,\Delta_2$:
\begin{enumerate}
\item Compute $D_0=B_1B_2$, $D_1=A_1B_2+A_2B_1$, and $D_2=A_1A_2$ modulo $q_l$. Their phase is the product of the two input phases.
\item Relinearize using either of the preceding methods, obtaining $(A,B)=(D_1,D_0)+C_\beta$ with phase $\Delta_1\Delta_2 M_1M_2+F$; $F$ includes multiplication and key-switching errors.
\item If the schedule calls for rescaling, set $q_{l-1}=q_l/w_l$ and
\[
(\widehat A,\widehat B)=[(\lceil A/w_l\rfloor,\lceil B/w_l\rfloor)]_{q_{l-1}},
\quad \Delta_{\mathrm{out}}=\Delta_1\Delta_2/w_l.
\]
The output error is $F/w_l+R_B+R_AS$, subject to the no-wrap and precision bounds above.
\end{enumerate}
\para{Swapping the Order of \textsf{Relinearization} and \textsf{Rescaling}:} It is possible to rescale the three-component ciphertext first, using rounding error $R_0+R_1S+R_2S^2$, and then relinearize with a key valid at the lower modulus. The two schedules are not identical ciphertext computations and have different error and cost bounds.
\end{tcolorbox}


\subsection{Ciphertext-to-Plaintext Multiplication}
\label{subsec:ckks-mult-plain}

Encode a plaintext multiplier $\Lambda$ at a chosen scale $\Delta_\Lambda$, giving $P_\Lambda=\lceil\Delta_\Lambda\Lambda\rfloor$. For a ciphertext $(A,B)$ at scale $\Delta$, compute
\[
(A',B')=(AP_\Lambda,BP_\Lambda)\pmod q,
\qquad \Delta'=\Delta\Delta_\Lambda.
\]
The phase includes both the input ciphertext error multiplied by $P_\Lambda$ and the plaintext encoding error. If rescaling by $w_l$ is needed, use the procedure in \autoref{subsubsec:ckks-mult-cipher-summary}, with output scale $\Delta\Delta_\Lambda/w_l$ and its rounding error. This common schedule consumes one level. It is not mandatory for every plaintext multiplication: an exactly represented small integer constant may use $\Delta_\Lambda=1$, and rescaling can sometimes be delayed when the modulus and precision budgets allow.


\subsection{\textsf{ModDrop}}
\label{subsec:ckks-moddrop}

For $q_{l-1}\mid q_l$, define
\[
(A',B')=([A]_{q_{l-1}},[B]_{q_{l-1}}).
\]
If $A'=A-K_Aq_{l-1}$ and $B'=B-K_Bq_{l-1}$, then
\[
A'S+B'=AS+B-(K_AS+K_B)q_{l-1}.
\]
Consequently the new phase is congruent to the old phase modulo $q_{l-1}$. If $\|\Delta M+E\|_\infty<q_{l-1}/2$, its centered representative remains exactly $\Delta M+E$; otherwise the phase can wrap and change the decoded value.

\begin{tcolorbox}[title={\textbf{\tboxlabel{\ref*{subsec:ckks-moddrop}} CKKS \textsf{ModDrop}}}]
Reduce each ciphertext component modulo a divisor of its current modulus. The tracked scale and the existing phase error are unchanged only when the complete phase fits the new centered interval. The modulus headroom decreases.
\end{tcolorbox}

\subsubsection{Difference between Modulus Switch and \textsf{ModDrop}}
\label{subsubsec:ckks-moddrop-vs-modswitch}

Both operations lower the ciphertext modulus. Rescaling also divides the tracked scale and old error by $q_l/q_{l-1}$ and introduces rounding error. \textsf{ModDrop} leaves the scale unchanged and preserves the phase only under its no-wrap condition. It is also used to align ciphertexts at different levels before addition or multiplication, not only when no future multiplication remains.

\subsection{Homomorphic Key Switching}
\label{subsec:ckks-key-switching}

Choose an exact integer gadget $g_i$ as in \autoref{subsubsec:relinearization-gadget-decomposition}, and decompose $A\equiv\sum_i d_ig_i\pmod q$. Evaluation keys $(a_i,b_i)$ under the new secret $S'$ satisfy $b_i+a_iS'\equiv g_iS+E_i\pmod q$. Then
\begin{tcolorbox}[title={\textbf{\tboxlabel{\ref*{subsec:ckks-key-switching}} CKKS Key Switching}}]
\[
(A',B')=(0,B)+\sum_i d_i(a_i,b_i),\qquad
B'+A'S'\equiv B+AS+\sum_i d_iE_i\pmod q.
\]
\end{tcolorbox}
The scale is unchanged and the error increases by the stated key-switching term. Correct interpretation requires the new phase to fit the modulus. This is the same underlying gadget technique used in \autoref{subsec:bfv-key-switching}, with CKKS's approximate decoding semantics.


\subsection{Homomorphic Rotation of Input Vector Slots}
\label{subsec:ckks-rotation}

The same root-exponent permutation underlies BFV and CKKS slot rotation, but BFV evaluates at roots in its plaintext ring and CKKS at complex roots. In the encoding convention of \autoref{subsec:ckks-encoding-decoding}, the inverse transform is $n^{-1}\hathat W I_n^R$, and the decoding transform is $\hathat W^*$. The matrix order matters.

For $J(j)=5^j\bmod2n$, the automorphism $M(X)\mapsto M(X^{J(h)})$ sends slot $j$ to the old slot $j+h\bmod(n/2)$, because $5^j5^h=5^{j+h}$. With $J_*(j)=-5^j\bmod2n$, the conjugate half rotates in the same direction.
\begin{tcolorbox}[title={\textbf{\tboxlabel{\ref*{subsec:ckks-rotation}} CKKS Homomorphic Rotation of Input Vector Slots}}]
\begin{enumerate}
\item Apply $X\mapsto X^{J(h)}$ to both $A$ and $B$, reducing modulo $X^n+1$. The secret becomes $S(X^{J(h)})$.
\item Use the exact gadget key switching in \autoref{subsec:ckks-key-switching} from $S(X^{J(h)})$ to $S(X)$. Include the added key-switching error in the output phase.
\end{enumerate}
The result rotates the $n/2$ independent complex slots left by $h$ positions.
\end{tcolorbox}

\para{Rotation within Half Slots:} A single $5^h$ automorphism rotates each of the two orbits separately. This does not prohibit other permutations constructed using masks and several automorphisms. CKKS's second half is determined by conjugation, while fully split BFV batching allows all $n$ slots to be independently chosen.


\subsubsection{Example}
\label{subsubsec:ckks-rotation-ex}

In this subsection, we will show the following 2 examples:
\begin{enumerate}
\item Encode an input vector $\vec{v}$ into a plaintext polynomial $M(X)$ based on our updated updated encoding \& decoding matrices $\hathat W$ and $\hathat W^*$
\item Rotate all elements of the input vector $\vec{v}$ $h$ positions to the left by updating the encoded plaintext $M(X)$ to $M(X^{J(h)})$
\end{enumerate}

$ $

We will use the same example of the input vector $\vec{v}$ used in \autoref{subsubsec:ckks-encoding-ex}: $\vec{v} = (1.1 + 4.3i, 3.5 - 1.4i)$.

Remember that the encoded plaintext polynomial of $\vec{v}$ is as follows:

$\Delta M(X) = 2355 + 1195X + 1485X^2 + 2933X^3 \in \mathcal{R}_{\langle 4 \rangle}=\mathbb{R}[X]/(X^4+1)$


Suppose we want to rotate the input vector $\vec{v}$ by 1 position to the left as follows:

$\vec{v}^{\langle h=1 \rangle} = (3.5 - 1.4i, 1.1 + 4.3i)$

$ $

Therefore, we update $\Delta M(X)$ to $\Delta M(X^{J(1)})$ as follows:

$\Delta M(X^{J(1)}) = \Delta M(X^{5}) = 2355 + 1195(X^5) + 1485(X^5)^2 + 2933(X^5)^3$

$= 2355 + 1195X^5 + 1485X^{10} + 2933X^{15}$

$= 2355 + 1195X\cdot(-1) + 1485X^{2}\cdot(-1)\cdot(-1) + 2933X^{3}\cdot(-1)\cdot(-1)\cdot(-1)$

$= 2355 - 1195X + 1485X^2 - 2933X^3$

$ $

The rotated \textit{forward-ordered} Hermitian input vector is computed as follows:

$\dfrac{\hathat W^* \cdot \Delta \vec{m}}{\Delta} = \begin{bmatrix}
1,\omega,\omega^2, \omega^3\\
1,\omega^5, \omega^{10},\omega^{15}\\
1,\overline{\omega}, \overline{\omega^2}, \overline{\omega^3}\\
1,\overline{\omega^5}, \overline{\omega^{10}}, \overline{\omega^{15}}
\end{bmatrix} \cdot \begin{bmatrix}
2355\\- 1195\\1485\\-2933
\end{bmatrix}\cdot \dfrac{1}{1024}$

$= \dfrac{\hathat W^* \cdot \Delta \vec{m}}{\Delta} = \begin{bmatrix}
1,\omega,\omega^2, \omega^3\\
1,\omega^5, \omega^{2},\omega^{7}\\
1,\overline{\omega}, \overline{\omega^2}, \overline{\omega^3}\\
1,\overline{\omega^5}, \overline{\omega^{2}}, \overline{\omega^{7}}
\end{bmatrix} \cdot \begin{bmatrix}
2355\\- 1195\\1485\\-2933
\end{bmatrix}\cdot \dfrac{1}{1024}$

$= \dfrac{\hathat W^* \cdot \Delta \vec{m}}{\Delta} = \begin{bmatrix}
1,\omega,\omega^2, \omega^3\\
1,\omega^5, \omega^{2},\omega^{7}\\
1,\omega^7, \omega^6, \omega^5\\
1,\omega^3, \omega^{6}, \omega
\end{bmatrix} \cdot \begin{bmatrix}
2355\\- 1195\\1485\\-2933
\end{bmatrix}\cdot \dfrac{1}{1024}$

$=
\begin{bmatrix}
1,\dfrac{\sqrt{2}}{2} + \dfrac{i\sqrt{2}}{2},i, -\dfrac{\sqrt{2}}{2} + \dfrac{i\sqrt{2}}{2}\\
1,-\dfrac{\sqrt{2}}{2} - \dfrac{i\sqrt{2}}{2}, i,\dfrac{\sqrt{2}}{2} - \dfrac{i\sqrt{2}}{2}\\
1,\dfrac{\sqrt{2}}{2} - \dfrac{i\sqrt{2}}{2}, -i, -\dfrac{\sqrt{2}}{2} - \dfrac{i\sqrt{2}}{2}\\
1,-\dfrac{\sqrt{2}}{2} + \dfrac{i\sqrt{2}}{2}, -i, \dfrac{\sqrt{2}}{2} + \dfrac{i\sqrt{2}}{2}
\end{bmatrix}$
$\cdot \begin{bmatrix}
2.2998046875\\-1.1669921875\\1.4501953125\\-2.8642578125
\end{bmatrix}$

$ $

$ \approx (3.500 - 1.4003i, \text{ } 1.0997 + 4.3007i,\text{ } 3.500 + 1.4003i, \text{ } 1.0997 - 4.3007i)$

$ $

Extract the first $\dfrac{n}{2} = 2$ elements in the above Hermitian vector to recover the input vector:

$(3.500 - 1.4003i, \text{ } 1.0997 + 4.3007i)$

$ \approx (3.5 - 1.4i, \text{ } 1.1 + 4.3i)$
$ = \vec{v}^{\langle h=1 \rangle}$ \textcolor{red}{  $\rhd$ The original input vector $\vec{v}$ rotated by 1 position to the left}



$ $

In an actual evaluation, apply the automorphism to the ciphertext components, obtaining
\[
(A(X^{J(1)}),B(X^{J(1)})).
\]
This ciphertext has secret $S(X^{J(1)})$. Key-switch it back to $S(X)$ before decrypting with the original key. The decoded result approximates the rotated vector $\vec v^{\langle h=1\rangle}$, with the original encoding error and the additional key-switching error.

$ $


\para{Source Code:} Examples of CKKS's homomorphic input vector slot rotation can be executed by running \href{https://github.com/wang2000star/fhe-textbook-cn/blob/chinese-edition/source%20code/ckks.py}{\underline{this Python script}}.





\subsection{Contemplation on CKKS Encoding}
\label{subsec:ckks-encoding-contemplate}

\para{Why CKKS's Encoding Uses the $\bm{(\mu=2n)}$-th Cyclotomic Polynomial:} Evaluation at the $n$ distinct roots of $X^n+1$ gives an invertible Vandermonde matrix. Any distinct evaluation points give an invertible Vandermonde matrix over a field; roots of unity additionally make evaluation compatible with multiplication modulo $X^n+1$, provide conjugate pairing, and allow efficient Fourier transforms. Power-of-two cyclotomic rings are a useful implementation choice, not a consequence of a general claim that only their Vandermonde matrices are invertible.

For the slot order $\omega^{5^j}$, applying $X\mapsto X^{5^h}$ permutes the rows of the evaluation (decoding) matrix. The inverse transform is the encoding matrix; the two matrix orientations must not be interchanged. Conjugate pairing reduces the independent complex coordinates to $n/2$ while keeping $n$ real coefficient coordinates.


\subsection{Homomorphic Conjugation}
\label{subsec:ckks-conjugation}

For real-coefficient $M$ and $\zeta_j=\omega^{J(j)}$, conjugation satisfies
\[
\overline{M(\zeta_j)}=M(\overline{\zeta_j})=M(\zeta_j^{-1}).
\]
The forward-ordered full vector therefore changes from $(\vec v,\overline{\vec v})$ to $(\overline{\vec v},\vec v)$. Use the same selected roots $J(j)=5^j\bmod 2n$ as in encoding, rather than replacing their order with $1,3,\ldots,n-1$.

\begin{tcolorbox}[title={\textbf{\tboxlabel{\ref*{subsec:ckks-conjugation}} CKKS Homomorphic Conjugation}}]
\begin{enumerate}
\item Apply the ring automorphism $X\mapsto X^{-1}=X^{2n-1}$ to both ciphertext components, with reduction modulo $X^n+1$.
\item Key-switch from $S(X^{-1})$ to $S(X)$, accounting for the added key-switching error.
\end{enumerate}
\end{tcolorbox}

\subsection{Sparsely Packing Ciphertexts}
\label{subsec:ckks-sparse-packing}

In \autoref{subsec:ckks-encoding-decoding}, we learned the CKKS encoding scheme, which encodes an input vector with $\dfrac{n}{2}$ slots (i.e., $\dfrac{n}{2}$-slot input vector) into an $(n-1)$-degree polynomial. While the polynomial ring's degree $n$ is fixed at the setup stage of CKKS as a security parameter, some applications may only need to use fewer than $\dfrac{n}{2}$ input vector slots. Suppose we only need to use $\frac{n'}{2}$ slots out of $\frac{n}{2}$ slots, where $n'\geq2$ divides $n$. Then, the corresponding input vector and encoded polynomial acquire a special property as described below:


\begin{tcolorbox}[title={\textbf{\tboxlabel{\ref*{subsec:ckks-sparse-packing}} CKKS's Sparsely Packing Polynomial and Ciphertext}}]

Suppose that an $\frac{n}{2}$-slot input vector gets encoded into a polynomial in $\mathbb{R}[X] / (X^n + 1)$. And suppose that $n'\geq2$ divides $n$. We define polynomial $M(X) \in \mathbb{R}[X] / (X^n + 1)$ as the one which has non-zero constants at the terms whose power is a multiple of $\frac{n}{n'}$ and all other terms have zero constants (i.e., $M(X) = c_0 + c_{\frac{n}{n'}}X^{\frac{n}{n'}} + c_{\frac{2n}{n'}}X^{\frac{2n}{n'}} + \cdots + c_{n - \frac{n}{n'}}X^{n - \frac{n}{n'}}$). We can express $M(X)$ as some $M_Y(Y) \in \mathbb{R}[Y] / (Y^{n'} + 1)$ where $Y = X^{\frac{n}{n'}}$ and thus $M(X) = M_Y(X^{\frac{n}{n'}})$.

$ $

Then, the following are true:

\begin{enumerate}
    \item Every $M_Y(Y) \in \mathbb{R}[Y]/(Y^{n'} + 1)$ is isomorphically mapped to (i.e., decoded into) some $\dfrac{n}{2}$-slot input vector which comprises $\dfrac{n}{n'}$ repetitions of $\dfrac{n'}{2}$ consecutive slot values.
    \item Conversely, if an $\dfrac{n}{2}$-slot input vector comprises $\dfrac{n}{n'}$ repetitions of the first $\dfrac{n'}{2}$ consecutive slot values, then the vector gets encoded into some $M_Y(Y) \in \mathbb{R}[Y] / (Y^{n'} + 1)$ (i.e., some polynomial in $\mathbb{R}[X] / (X^n + 1)$ that has zero constants at the terms whose degree is not a multiple of $\dfrac{n}{n'}$).
\end{enumerate}

\end{tcolorbox}

Assume $n'\geq2$ divides the power-of-two $n$. The forward inclusion follows below. To infer the converse, also note that both the subring $\mathbb{R}[Y]/(Y^{n'}+1)$ and the space of repeated slot vectors have real dimension $n'$. The injective canonical embedding maps the first into the second, so equality of dimensions proves surjectivity onto that subspace. An isomorphism alone, without this subspace argument, would not establish the converse.

\subsubsection{Forward Proof: Decoding of Sparsely Packed Ciphertext}

We will prove that for each $M_Y(Y) \in \mathbb{R}[Y] / (Y^{n'} + 1)$ (i.e., a polynomial in $\mathbb{R}[X] / (X^n + 1)$ that has non-zero constants only at those terms with a power that is a multiple of $\dfrac{n}{n'}$ and zero constants in all other terms), the polynomial is decoded into some $\dfrac{n}{2}$-slot input vector which comprises $\dfrac{n}{n'}$ repetitions of the first $\dfrac{n'}{2}$ consecutive slot values.

To decode $M(X)$ into an input vector, evaluate $M(X)$ at the selected $\frac{n}{2}$ roots among the $n$ distinct primitive $(\mu=2n)$-th roots of $X^n + 1$. The evaluation vector is:

$\bm{(} \text{ }
M(\omega^{J(0)}), \text{ } M(\omega^{J(1)}), \cdots,  M(\omega^{J(\frac{n}{2}-1)})\bm{)}$

\textcolor{red}{ $\rhd$ where $\omega=e^{\frac{i\pi}{n}}$, the base and generator of the primitive ($\mu=2n$)-th roots of unity}

$ $

But since $M(X) = M_Y(X^{\frac{n}{n'}})$, the above evaluation is equivalent to evaluating:

$\bm{(} \text{ }
M_Y((\omega^{J(0)})^{\frac{n}{n'}}), \text{ } M_Y((\omega^{J(1)})^{\frac{n}{n'}}), \cdots,  M_Y((\omega^{J(\frac{n}{2}-1)})^{\frac{n}{n'}}) \bm{)}$

$ $

$= \bm{(} \text{ }
M_Y((\omega^{\frac{n}{n'}})^{J(0)}), \text{ } M_Y((\omega^{\frac{n}{n'}})^{J(1)}), \cdots,  M_Y((\omega^{\frac{n}{n'}})^{J(\frac{n}{2}-1)})\bm{)}$

$ $

$= \bm{(} \text{ }
M_Y(\xi^{J(0)}), \text{ } M_Y(\xi^{J(1)}), \cdots,  M_Y(\xi^{J(\frac{n}{2}-1)})\bm{)}$

\textcolor{red}{ $\rhd$ where $\xi = e^{\frac{i\pi}{n'}}$, the base and generator of the primitive $(\mu=2n')$-th roots of unity}

$ $

Notice that $\xi = \omega^{\frac{n}{n'}}$. Therefore, the above evaluation of $M_Y(Y)$ outputs $\dfrac{n}{n'}$ repeated values of $M_Y(Y)$ evaluated at $\frac{n'}{2}$ distinct primitive $(\mu=2n')$-th roots of unity.



\subsection{Modulus Bootstrapping}
\label{subsec:ckks-bootstrapping}

\textbf{- Reference:} \href{https://eprint.iacr.org/2018/153.pdf}{Bootstrapping for Approximate Homomorphic Encryption}~\cite{ckks}

In a standard leveled evaluation schedule, rescaling eventually exhausts the available modulus chain. Bootstrapping starts from a small modulus $q_0$, raises it to a working modulus $q_L$, and returns a ciphertext at some $q_l>q_0$ after consuming levels internally. The output need not be at the maximum modulus $q_L$.

Take centered representatives of $(A,B)\in R_{q_0}^2$ and suppose $[B+AS]_{q_0}=\Delta M+E$. Their unreduced phase is
\[
B+AS=\Delta M+E+q_0K,
\]
where $K\in R$. This is an integer-lift identity, not a decryption identity with the modulo operation already performed. Bootstrapping approximately removes the $q_0K$ term while retaining the previous numerical error and introducing a new error.

\subsubsection{High-level Idea}
\label{subsubsec:ckks-bootstrapping-high-level}

As the first step of bootstrapping, we forcibly change the modulus of the ciphertext $(A, B)$ from $q_0$ to $q_L$. Then, its decryption with reduction modulo $q_L$ would output:

$\textsf{RLWE}^{-1}\textbf{(}\textsf{ct} = (A, B)\textbf{)} \bmod q_L$

$ = B + A\cdot S \bmod q_L$

$= \Delta M + E + q_0 K \bmod q_L$

This interpretation requires $\|B+AS\|_\infty<q_L/2$. For centered input representatives, $\|B+AS\|_\infty\leq(q_0/2)(1+\|S\|_1)$, so a sufficiently large $q_L$ can ensure it. Ternary coefficients alone are not a dimension-independent bound; $\|S\|_1$ can be as large as $n$.


In the $\Delta M + E + q_0 K \bmod q_L$ term, notice that because of the $q_0K$ term which is not modulo-reduced by $q_0$ anymore, the ciphertext's decrypted plaintext polynomial's each $i$-th term would get a corrupted coefficient $\Delta m_i + e_i + q_0\cdot k_i \bmod q_L$ instead of $\Delta m_i + e_i \bmod q_L$. So, we now need to eliminate the garbage term $q_0\cdot k_i \bmod q_L$ in each coefficient and distill the pure plaintext coefficient $\Delta m_i + e_i$.



\begin{figure}[h!]
    \centering
  \includegraphics[width=1.0\linewidth]{figures/modulo-reduction-sine.png}
  \caption{Sine function $f(x) = \dfrac{q_0}{2\pi}\cdot \sin \left(\dfrac{2\pi x}{q_0}\right)$ such that $f(\Delta m_i + e_i + q_0k_i) \approx \Delta m_i + e_i$ (provided $\Delta m_i + e_i \ll q_0$) \href{https://eprint.iacr.org/2018/153.pdf}{(Source)}}
  \label{fig:modulo-reduction-sine}
\end{figure}

The function is periodic with period $q_0$. Near $x=kq_0$, it approximates the local centered remainder $x-kq_0$, not the global function $x$. Its derivative at $kq_0$ is $1$.

$ $

Combining these two properties, given input $x = \Delta m_i + e_i + q_0k_i$,

\[
\begin{aligned}
f(\Delta m_i+e_i+q_0k_i)
&=\frac{q_0}{2\pi}\sin\left(\frac{2\pi(\Delta m_i+e_i+q_0k_i)}{q_0}\right)\\
&=\frac{q_0}{2\pi}\sin\left(\frac{2\pi(\Delta m_i+e_i)}{q_0}\right)\\
&\approx\Delta m_i+e_i.
\end{aligned}
\]

$ $

, provided $|\Delta m_i+e_i|/q_0$ is sufficiently small for the target accuracy. This is an additional bootstrapping parameter condition, not a consequence of ordinary CKKS decryption correctness. For $r_i=\Delta m_i+e_i$, the intrinsic sine error satisfies
\[
\left|\frac{q_0}{2\pi}\sin(2\pi r_i/q_0)-r_i\right|
\leq\frac{2\pi^2}{3q_0^2}|r_i|^3.
\]
Polynomial approximation and homomorphic evaluation add further errors. Since ciphertext multiplication acts slot-wise, the algorithm first applies a linear coefficient-to-slot transform, evaluates the nonlinear approximation in slots, and then applies the inverse linear transform, as outlined in \autoref{tab:ckks-bootstrapping-procedure}.

\begin{table}
\begin{tabular}{|lp{0.78\linewidth}|}
\hline
1 & \textsf{ModRaise}: lift ciphertext coefficients to $q_L$, obtaining phase $\Delta M+E+q_0K$.\\
2 & \textsf{CoeffToSlot}: homomorphically apply the inverse embedding to put the $n$ real phase coefficients in two ciphertexts.\\
3 & \textsf{EvalExp}: evaluate a nonlinear polynomial approximation to the periodic modular-reduction function, with an explicit input range and error budget.\\
4 & \textsf{SlotToCoeff}: repack the approximate coefficients into one ciphertext at the remaining modulus $q_l>q_0$, including all accumulated errors.\\
\hline
\end{tabular}
\caption{High-level Description of CKKS's Bootstrapping Procedure}
\label{tab:ckks-bootstrapping-procedure}
\end{table}

\subsubsection{Mathematical Expansion of the High-level Idea}
\label{subsubsec:ckks-bootstrapping-high-level-correctness}

The following identities describe plaintext values represented by ciphertexts; they are not operations that expose the secret key or decrypt during evaluation. Let $T=\Delta M+E+q_0K$ be the lifted phase, and use $\mathsf{Enc}(T)$ only as shorthand for a ciphertext representing that phase.

\para{\textsf{\textbf{\underline{ModRaise}:}}} Lift the coefficients of $(A,B)$ to the working modulus $q_L$. Under the bound above, the lifted ciphertext has raw phase $T$. It is a valid starting point for homomorphic evaluation without computing $AS$ in the clear.

\para{\textsf{\textbf{\underline{CoeffToSlot}:}}} Let $\sigma$ be the canonical embedding and let $\pi$ select the $n/2$ independent slots. The real-linear inverse embedding maps the slots $\pi\sigma(T)$ to the coefficient vector $(T_i)_{i=0}^{n-1}$. Its homomorphic implementation uses rotations, conjugation, additions, and multiplication by encoded matrix constants. Store the $n$ real coefficients in two ciphertexts, with $n/2$ coefficients per ciphertext. Scales and the errors of these transformations must be tracked.

\para{\textsf{\textbf{\underline{EvalExp}:}}} Homomorphically evaluate a polynomial circuit $\widehat f$ approximating $f(x)=q_0\sin(2\pi x/q_0)/(2\pi)$ on the required input range. This is a nonlinear computation, not a linear transformation or multiplication by one encoded polynomial. For each coefficient,
\[
\widehat f(T_i)=\Delta m_i+e_i+e_{b,i},
\]
where $e_{b,i}$ includes sine error, polynomial-approximation error, and all evaluation errors.

\para{\textsf{\textbf{\underline{SlotToCoeff}:}}} Apply the inverse linear packing operation to the two resulting ciphertexts. The final phase is represented as $\Delta M+E+E_b$ at its tracked output scale and modulus $q_l$; $E_b$ also includes the final transform errors.

\para{\textsf{\textbf{\underline{Bootstrapping Without a Secret Key}:}}} The actual circuit operates on the lifted ciphertext throughout. It does not decrypt and re-encrypt, nor does an identity such as $\mathsf{Enc}(\mathsf{Dec}(\textsf{ct}))=\textsf{ct}$ hold for randomized noisy encryption. Public evaluation keys permit the rotations, conjugations, and relinearizations. The output represents an approximation to the original decoded value, with a larger available modulus and additional error.

\subsubsection{Details: \textsf{CoeffToSlot}}
\label{subsubsec:ckks-bootstrapping-coefftoslot-details}

Set $\zeta_j=\omega^{J(j)}$ and define the $n/2\times n$ evaluation matrix $U$ by $U_{j,k}=\zeta_j^k$. If $\vec t$ is the real coefficient vector of $T$, its independent slots are $\vec z=U\vec t$. Orthogonality of the roots gives
\[
\vec t=\frac{1}{n}(\overline U^{\,T}\vec z+U^T\overline{\vec z}).
\]
Split $U=(U_0\;U_1)$ into two square blocks, according to the first and second halves of the coefficient indices. Then
\[
\vec t_0=\frac{1}{n}(\overline{U_0}^{\,T}\vec z+U_0^T\overline{\vec z}),\qquad
\vec t_1=\frac{1}{n}(\overline{U_1}^{\,T}\vec z+U_1^T\overline{\vec z}).
\]
Applying these two real-linear maps homomorphically produces ciphertexts whose slots contain the two coefficient halves, up to encoding and evaluation errors. Conjugation is computed as in \autoref{subsec:ckks-conjugation}; matrix-vector products use encoded diagonals and rotations as in \autoref{subsec:bfv-matrix-multiplication}. They are ciphertext--plaintext multiplications and additions, with level consumption determined by the chosen scale schedule. The notation does not mean multiplying a matrix directly by a ciphertext polynomial's coefficient list in the clear.

\subsubsection{Details: \textsf{EvalExp}}
\label{subsubsec:ckks-bootstrapping-evalexp-details}

We will use the sine function $f(x) = \dfrac{q_0}{2\pi}\cdot \sin\left(\dfrac{2\pi x}{q_0}\right)$ to approximately eliminate $q_0k_i$ from $\Delta m_i + e_i + q_0k_i$ by computing $f(\Delta m_i + e_i + q_0k_i) \approx \Delta m_i + e_i$. This approximation works if $\Delta m_i + e_i$ is very close to $x=0$ relative to $q_0$ (i.e., $\Delta m_i + e_i \ll q_0$). Still, the elimination of $q_0k_i$ is approximate (i.e., $\approx \Delta m_i + e_i$), because $f(x)$ is $y \approx x$ nearby $x=0$, not exactly $y = x$.

One issue is that we need to evaluate $f(x)$ homomorphically based on \textsf{ct\textsubscript{s1}} and \textsf{ct\textsubscript{s2}} as inputs (i.e., $f(\textsf{ct\textsubscript{s1}})$ and $f(\textsf{ct\textsubscript{s2}})$), but FHE supports only $(+, \cdot)$ operations, whereas the sine graph cannot be formulated by only $(+, \cdot)$. Therefore, we will approximate the sine function $f(x)$ by using the Taylor series (\autoref{sec:taylor-series}):

$f(x) = f(a) + \dfrac{f'(a)}{1!}(x-a) + \dfrac{f''(a)}{2!}(x-a)^2 + \dfrac{f'''(a)}{3!}(x-a)^3 + \cdots = \sum\limits_{d=0}^{\infty}\dfrac{f^{(d)}(a)}{d!}(x-a)^d$

$ $

If we approximate $f(x)$ around $x = 0$, then the approximated polynomial is as follows:

$f(x) = \dfrac{q_0}{2\pi}\cdot \sin\left(\dfrac{2\pi}{q_0}\cdot 0\right) + \dfrac{q_0}{2\pi}\cdot\dfrac{2\pi}{q_0}\cdot\dfrac{\cos\left(\dfrac{2\pi}{q_0} \cdot 0\right)}{1!}\cdot x + \dfrac{q_0}{2\pi}\cdot \left(\dfrac{2\pi}{q_0}\right)^{2}\cdot\dfrac{-\sin\left(\dfrac{2\pi}{q_0} \cdot 0\right)}{2!}\cdot x^2 + \cdots$

$ = \dfrac{q_0}{2\pi}\cdot \sum\limits_{j=0}^{\infty}\left(\dfrac{(-1)^j}{(2j + 1)!}\cdot \left(\dfrac{2\pi x}{q_0}\right)^{2j + 1}\right)$

$ \approx \dfrac{q_0}{2\pi}\cdot \sum\limits_{j=0}^{h}\left(\dfrac{(-1)^j}{(2j + 1)!}\cdot \left(\dfrac{2\pi x}{q_0}\right)^{2j + 1}\right) = \hat f(x)$

$ $

, where $\hat{f}(x)$ is a $(2h + 1)$-degree polynomial.

For centered representatives of $A,B$ and a small secret $S$, the integer phase satisfies $B+AS=\Delta M+E+q_0K$. Its coefficient norm is at most $(q_0/2)(1+\|S\|_1)$, and the bound on $K$ also includes $\|\Delta M+E\|_\infty/q_0$. Choose $k_{\max}$ to bound the entire normalized input $|x|/q_0$, including the residual message and error, not just $|k_i|$. Sparse secrets can reduce this range, but require a security analysis with their actual distribution.

The normalized approximation interval is controlled by $\theta=2\pi x/q_0$, so increasing $q_0$ alone does not increase its width when $|x|/q_0$ is fixed. For a degree-$d$ polynomial, Horner evaluation uses $d$ multiplications, and standard parallel power constructions have logarithmic multiplicative depth with polynomial work. The cost is not generally exponential in the degree. Scaling and repeated squaring exploit the special exponential-function structure to reduce the required work and depth:

\[
\begin{cases}
e^{i\cdot \theta} = \cos \theta + i \cdot \sin \theta\\
(e^{i \cdot \theta})^2 = e^{i\cdot 2\theta}\\
\end{cases}
\]

Set $\theta=2\pi x/q_0$. First approximate $e^{i\theta/2^r}$ on the smaller interval $|\theta|/2^r\leq2\pi k_{\max}/2^r$, then square the result $r$ times to approximate $e^{i\theta}$. The required Taylor degree depends on this normalized interval and the desired final error. Squaring amplifies the initial approximation error and adds arithmetic error, so the initial tolerance must account for all $r$ squarings.

Using the Taylor series with degree $d_0$ around $x = 0$, we can approximate $e^{\frac{2\pi i x}{2^r q_0}}$ as:

$f_e(x) = e^{\frac{2\pi i x}{2^r q_0}} \approx \sum\limits_{d=0}^{d_0}\dfrac{1}{d!}\left(\dfrac{2\pi i x}{2^rq_0}\right)^d = \hat{f}_e(x)$

$ $

Then, we iteratively square $\hat{f}_e$ total $r$ times to get:

$(\hat{f}_e(x))^{2^r} \approx (f_e(x))^{2^r} = e^{i\frac{2\pi x}{q_0}} = e^{i\theta}$

$ $

Then, based on Euler's formula $e^{i\cdot \theta} = \cos \theta + i \cdot \sin \theta$, we can derive the following relations:

$ $

$\overline{e^{i\cdot \theta}} = \cos \theta + \overline{i \cdot \sin \theta} $

$ e^{-i\cdot \theta} = \cos \theta - i \cdot \sin \theta$

$ e^{i\cdot \theta} - e^{-i\cdot \theta} = (\cos \theta + i \cdot \sin \theta) - (\cos \theta - i \cdot \sin \theta) = 2i\sin \theta$

$\sin \theta = \dfrac{-i}{2} \cdot (e^{i\cdot \theta} - e^{-i\cdot \theta})$

$\dfrac{q_0}{2\pi}\cdot \sin \theta = \dfrac{q_0}{2\pi}\cdot \dfrac{-i}{2} \cdot (e^{i\cdot \theta} - e^{-i\cdot \theta})$

$ $

Substituting $\theta = \dfrac{2\pi x}{q_0}$, we finally get:

$\dfrac{q_0}{2\pi}\cdot \sin \left(\dfrac{2\pi x}{q_0}\right) = \dfrac{q_0}{2\pi}\cdot \dfrac{-i}{2} \cdot (e^{i\cdot \frac{2\pi x}{q_0}} - e^{-i\cdot \frac{2\pi x}{q_0}})$

$ $

Using the final relation above, the \textsf{EvalExp} step homomorphically evaluates the approximation of $\dfrac{q_0}{2\pi}\cdot \sin \left(\dfrac{2\pi x}{q_0}\right)$ where $x = \Delta m_i + e_i + q_0k_i$ as follows:

\begin{enumerate}
\item Homomorphically evaluate the exponential approximation, obtaining $\widehat z\approx e^{2\pi ix/q_0}$ with a quantified error.

\item Apply homomorphic conjugation (\autoref{subsec:ckks-conjugation}) to obtain an approximation to $\overline{\widehat z}$, accounting for its key-switching error.

\item Subtract the two results and multiply by the encoded constant $-iq_0/(4\pi)$. This approximates $\frac{q_0}{2\pi}\sin(2\pi x/q_0)$. The factor $q_0/(2\pi)$ is essential; multiplying only by $-i/2$ would return the unscaled sine.
\end{enumerate}

$ $

The result of \textsf{EvalExp} is two ciphertexts whose input vector slots store the bootstrapped coefficients of $M_c$, which are modulo-reduced $q_0$ from $\Delta m_i + e_i + q_0k_i$ to $\Delta m_i + e_i + e_{bi}$. Note that $e_{bi}$ is a bootstrapping error introduced by the following three factors: (1) the intrinsic homomorphic $(+, \cdot)$ computation noises of the \textsf{CoeffToSlot}, \textsf{EvalExp}, and \textsf{SlotToCoeff} steps; (2) the \textsf{EvalExp} step's Taylor polynomial approximation error of the exponential function $e^{i\theta}$; (3) the \textsf{EvalExp} step's sine graph error, since the graph is not exactly $y=x$ around $x=0$, but only $y \approx x$.

Note that since the output of the \textsf{CoeffToSlot} step was split into 2 ciphertexts (\textsf{ct\textsubscript{s1}} and \textsf{ct\textsubscript{s2}}), the output of the \textsf{EvalExp} step is also in 2 ciphertexts: (\textsf{ct\textsubscript{b1}} and \textsf{ct\textsubscript{b2}}). The input vector slots of \textsf{ct\textsubscript{b1}} store $(\Delta m_i + e_i + e_{bi})_{i=0}^{\frac{n}{2} - 1}$, whereas the input vector slots of \textsf{ct\textsubscript{b2}} store $(\Delta m_i + e_i + e_{bi})_{i=\frac{n}{2}}^{n - 1}$.


\subsubsection{Details: \textsf{SlotToCoeff}}
\label{subsubsec:ckks-bootstrapping-slottocoeff-details}

Using the same blocks $U_0,U_1$ as in \autoref{subsubsec:ckks-bootstrapping-coefftoslot-details}, let the two EvalExp output vectors be $\vec b_0,\vec b_1$, containing approximations to the first and second coefficient halves. Compute
\[
\vec z_{\mathrm{out}}=U_0\vec b_0+U_1\vec b_1.
\]
The resulting single ciphertext has $n/2$ complex slots representing the polynomial whose $n$ real coefficients are the concatenation of those halves. Its implicit conjugate slots are $\overline{\vec z_{\mathrm{out}}}$; they are not a second encryption of the same coefficient polynomial. The homomorphic transform incurs encoding, key-switching, and scaling errors. Thus its raw output phase is $\Delta_{\mathrm{out}}M+E_{\mathrm{old,out}}+E_b$ at modulus $q_l$, with all scales and errors propagated through the chosen implementation. Under a schedule restoring the input scale, this is abbreviated $\Delta M+E+E_b$.

\subsubsection{Reducing the Bootstrapping Overhead by Sparsely Packing Ciphertext}
\label{subsubsec:ckks-bootstrapping-time-reduction}

Assume a proper even divisor $2\leq n'<n$ and set $d=n/n'$. Repeating $n'/2$ independent slots produces a message polynomial in $Y=X^d$, with $n'$ potentially nonzero coefficients. To replicate an initially zero-padded block, use $\log_2 d$ doubling rounds with shifts $-(n'/2)2^j$, for $0\leq j<\log_2d$.

Message sparsity alone is insufficient after modulus raising: the phase $\Delta M(Y)+E(X)+q_0K(X)$ generally has nonsparse $E$ and $K$. The original CKKS bootstrapping paper, Section 5.2, therefore inserts a partial-sum procedure after ModRaise and before CoeffToSlot:
\[
\textsf{ct}\leftarrow\textsf{ct}+\textsf{Rotate}(\textsf{ct},(n'/2)2^j),
\qquad 0\leq j<\log_2d.
\]
This trace-like operation annihilates coefficient positions not divisible by $d$ and multiplies retained coefficients by $d$, up to evaluation errors. Its normalization must be included in the tracked scales and the modular-reduction period. Omitting it does not justify discarding the other lifted coefficients.

The smaller canonical embedding uses the roots of $Y^{n'}+1$ and has $n'$ real coefficient coordinates, not $n/n'$. Since a proper power-of-two divisor satisfies $n'\leq n/2$, these coefficient values can fit in one ciphertext. The corresponding smaller coefficient/slot transforms and modular-reduction evaluation reduce cost. Their exact savings depend on the transform algorithm, rotations, and normalization; they are not automatically exactly one half for every parameter choice.

\subsubsection{Summary}
\label{subsubsec:ckks-bootstrapping-summary}

\begin{tcolorbox}[title={\textbf{\tboxlabel{\ref*{subsubsec:ckks-bootstrapping-summary}} CKKS Bootstrapping}}]
\begin{enumerate}
\item Start from $(A,B)\in R_{q_0}^2$ with centered phase $\Delta M+E$ and the required input-range bound.
\item Lift the ciphertext coefficients to $q_L$. The raw phase is $\Delta M+E+q_0K$, under a no-wrap bound at $q_L$. RNS implementations must compute any needed basis extension.
\item Apply CoeffToSlot, using two ciphertexts for a fully packed input, or the properly normalized sparse procedure above.
\item Evaluate the nonlinear polynomial approximation to modular reduction, budgeting both approximation error and homomorphic evaluation error.
\item Apply SlotToCoeff. The output is at a remaining modulus $q_l>q_0$, with tracked scale and accumulated error. In equal-scale shorthand, its phase is $\Delta M+E+E_b$, not exactly $\Delta M+E$.
\end{enumerate}
\textbf{Limitation:} Bootstrapping does not recover precision already lost in the input. Further numerical computation must remain within its message, error, and modulus budgets.
\end{tcolorbox}

\para{Comparison between BFV and CKKS Bootstrapping:} Exact-scheme bootstrapping can remove prior encryption noise by evaluating exact decryption, provided its correctness conditions hold. CKKS retains the prior numerical approximation error and adds a new error. Neither fixed-modulus BFV nor a fixed CKKS parameter set supports arbitrary unbounded work without satisfying the respective noise and accuracy conditions. The CKKS procedure here uses the lifted ciphertext directly; the absence of a separate encrypted-secret inner product is not explained by an assumption that all intermediate scales remain fixed.

\subsubsection{Reducing the Bootstrapping Noise}
\label{subsubsec:ckks-bootstrapping-noise-reduction}

As explained in \autoref{subsubsec:ckks-bootstrapping-evalexp-details}, the bootstrapping procedure generates three types of noises:

%\renewcommand{\labelenumi}{\underline{Type-\arabic{enumi} Noise:}}
\begin{itemize}
\item \underline{Type-1 Noise}: the intrinsic homomorphic $(+, \cdot)$ computation noises of the \textsf{CoeffToSlot}, \textsf{EvalExp}, and \textsf{SlotToCoeff} steps
\item \underline{Type-2 Noise}: the \textsf{EvalExp} step's approximation error of the exponential function $e^{i\theta}$
\item \underline{Type-3 Noise}: the \textsf{EvalExp} step's sine graph error (i.e., not exactly $y=x$ around $x=0$, but only $y \approx x$)
\end{itemize}


$ $

All three error sources must be bounded for the chosen input range and parameter set. Finite polynomial evaluation generally retains a nonzero approximation error, while its degree and precision can be increased to reduce it. Improved modular-reduction approximations can reduce the intrinsic sine error; they do not eliminate the need to analyze arithmetic and approximation errors together.

\begin{figure}[h!]
    \centering
  \includegraphics[width=1.0\linewidth]{figures/arc-sine.png}
  \caption{Arc-sine graph for smaller approximation error \href{https://www.google.com/search?sca_esv=c744cb070de47b7e&sxsrf=ADLYWIKKaF93fcAFvJC4ehJso1E5qFTaow:1732546684173&q=arcsin(sin(x))&source=lnms&fbs=AEQNm0DmKhoYsBCHazhZSCWuALW8l8eUs1i3TeMYPF4tXSfZ95GzcfXnm5XYTvJV_9Qreh2py964ICpZJthXkELijctC8pFBYULoa3-fvQmwK0VJF0ntzsbN_W2CCJL9N57SWFNwWI58jCKaBJSdPgkprHQVK8H1PYOYWXMHTCCV-rDbC44rR6ANM870jZCZRtTKwFWtMIe2&sa=X&ved=2ahUKEwjZy6qt3_eJAxXbhlYBHXD2IvEQ0pQJegQIDxAB&cshid=1732546705369143&biw=1280&bih=635&dpr=1.5}{(Source)}}
  \label{fig:arc-sine}
\end{figure}

$ $

\para{$\bm{\arcsin(\sin(x))}$ Approximation (\href{https://dl.acm.org/doi/10.1007/978-3-030-77870-5_22}{EUROCRYPT 2021}~\cite{cryptoeprint:2020/552}):} The exact composition equals $x$ only on a principal interval such as $[-\pi/2,\pi/2]$ and extends periodically with alternating slopes. Around each required multiple of the period it can represent the desired local remainder. Homomorphic implementations approximate the component functions on specified intervals, for example using minimax polynomials. Their approximation error, interval restrictions, and additional depth must be specified; a universal cost of exactly three levels does not follow from the identity.

$ $

\para{Meta-BTS (\href{https://dl.acm.org/doi/10.1145/3548606.3560696}{CCS 2022}~\cite{cryptoeprint:2022/1167}):} Meta-BTS improves precision by bootstrapping an amplified residual and subtracting it from the preceding result. The following conditional derivation makes the scale and input-range requirements explicit. It is a sufficient construction, not a parameter-independent claim about every implementation of the paper.

Write $v_D(c)$ for the canonical embedding of the centered decryption phase of $c$, divided by its scale $D$. All norms below are slot infinity norms unless stated otherwise. Let the base bootstrap take modulus $q$ and scale $\Delta$ to the remaining output modulus $Q$ and the same scale $\Delta$. Suppose its contract, for every admissible input with $\|v_\Delta(c)\|_\infty\leq B$, is
\[
v_\Delta(\mathsf{BTS}(c))=v_\Delta(c)+b(c),
\qquad \|b(c)\|_\infty\leq\varepsilon.
\]
Admissibility includes the internal lifting and approximation ranges, evaluation keys, and arithmetic bounds. Any output-scale adjustment is included in this contract and its error. Here $Q$ is the modulus remaining after the base bootstrap, not its temporary raised modulus. The error is measured relative to the actual decoded input, including any error that input already contains.

Assume a componentwise rescale ending at scale $D$ has decoded error at most $T/D$, and put $\rho=T/\Delta$. For nearest-integer component rounding in the degree-$n$ ring, a conservative deterministic choice is
\[
T=\frac{n}{2}(1+\|S\|_1),
\qquad
\|\sigma(R_B+R_AS)\|_\infty\leq T,
\]
where the rounding polynomials have coefficient magnitudes at most $1/2$, $\|S\|_1$ is the coefficient one-norm, and $\sigma$ is the canonical embedding. A tighter established bound may replace this one. The phase must remain within the centered interval at the destination modulus.

For two calls, take an actual integer rescale factor $a>1$, an input $c$ of modulus $aq$ and scale $a\Delta$, and $z_0=v_{a\Delta}(c)$. The following operations are componentwise ciphertext operations; brackets mean modulus reduction, without division or rounding:
\begin{align*}
u&=\mathsf{RS}_a(c), &v_\Delta(u)&=z_0+r,\\
y&=\mathsf{BTS}(u), &v_\Delta(y)&=z_0+r+b_1,\\
h&=ay, &v_{a\Delta}(h)&=z_0+r+b_1,\\
d&=[h]_q-[c]_q, &v_\Delta(d)&=a(r+b_1),\\
w&=\mathsf{BTS}(d), &v_\Delta(w)&=a(r+b_1)+b_2,\\
c_{\mathrm{out}}&=h-w, &v_{a\Delta}(c_{\mathrm{out}})&=z_0-b_2/a.
\end{align*}
The moduli are $q$ for $u,d$ and $Q$ for $y,h,w,c_{\mathrm{out}}$. In the last subtraction, the components of $w$ are interpreted at scale $a\Delta$; this intentional scale interpretation divides its decoded value by $a$. The first rescale error $r$ cancels together with $b_1$. The residual is bootstrapped with the sign ``new ciphertext minus input.'' In particular, ModDrop must not be replaced by a rescale.

Sufficient input-range conditions are
\[
\|z_0\|_\infty+\rho\leq B,
\qquad a(\rho+\varepsilon)\leq B.
\]
For example, $aq\mid Q$ and $q>2\Delta(B+\varepsilon)$ are conservative external modulus conditions ensuring the required reductions and centered phase bounds. The individual reductions of $h$ and $c$ to $q$ may wrap; their difference still has the stated small centered phase because reduction is a ring homomorphism and the residual bound holds. Under these conditions the added output error is at most $\varepsilon/a$. The identity holds for each actual execution, even if $r,b_1,b_2$ are correlated.

For $k$ calls, let $P_1=1$ and $P_{j+1}=a_jP_j$, where every $a_j>1$ is the actual available integer rescale factor. The $j$-call construction has input modulus $P_jq$ and input/output scale $P_j\Delta$. To construct the next one, rescale the new input by $a_j$, apply the $j$-call construction, multiply its components by $a_j$, and bootstrap the difference with the original input reduced to $q$, interpreting that difference at scale $\Delta$. If the first rescale error is $r_j$ and the preceding construction adds $e_j$, then the last base call receives
\[
P_{j+1}(r_j+e_j),\qquad
\|r_j\|_\infty\leq\rho/P_j,\qquad
\|e_j\|_\infty\leq\varepsilon/P_j.
\]
After subtracting its output, the resulting value is $z_0-b_{j+1}/P_{j+1}$. Induction therefore gives the sufficient conditions and bound
\begin{gather*}
a_j(\rho+\varepsilon)\leq B\quad(1\leq j<k),\qquad
\|z_0\|_\infty+\rho\sum_{j=1}^{k-1}P_j^{-1}\leq B,\\
P_kq\mid Q,\qquad q>2\Delta(B+\varepsilon),\qquad
\|v_{P_k\Delta}(c_{\mathrm{out}})-z_0\|_\infty\leq\varepsilon/P_k.
\end{gather*}
This proof uses $k$ base bootstrap calls and requires their contract on all the derived inputs. Its divisibility conditions are convenient sufficient conditions, not asserted to be necessary for every variant.

In the ideal notation $a_j=2^b$ and $\varepsilon=2^{-b}$, the error expression becomes $2^{-kb}$ and the scale is $2^{(k-1)b}\Delta$. However, with $B=1$ and any positive $\rho$, the bound $2^b(\rho+2^{-b})>1$ leaves no residual input margin. Thus a $b$-bit bound alone does not establish the advertised iteration: one needs a sharper error bound, a smaller gain, or a larger certified input range. In RNS, a factor close to $2^b$ must be tracked at its actual value in both $P_k$ and the scale. For instance, if the actual factor is $a+1$ but restoration uses $a$, the residual contains $-z_0+a(r+b_1)$ rather than only $a(r+b_1)$; even a small relative mismatch can violate the input range. Changing scale metadata alone does not recover precision.

\para{Limitation in Noise Handling:} A quantitative guarantee must combine the approximation domains, arithmetic errors, and the preceding residual-input conditions. Here is one deterministic error budget for the arcsine approach; it also explains which additional bounds a concrete implementation must supply. Let the ideal normalized coefficient values $t$ lie in
\[
D_{K,h}=\bigcup_{j=-K}^{K}[j-h,j+h],\qquad 0<h<1/4.
\]
On these real intervals, $\arcsin(\sin(2\pi t))/(2\pi)=t-j$. Let $p_s$ be the complete inner polynomial, including any cosine/double-angle composition, with uniform real-domain approximation error $\epsilon_s$. Let the CoeffToSlot error be at most $E_C$, and let $H_s$ bound the total inner circuit evaluation error at its output. Define
\[
L_s=\sup_{\operatorname{dist}(z,D_{K,h})\leq E_C}|p_s'(z)|,
\qquad E_s=L_sE_C+\epsilon_s+H_s,
\qquad s_h=\sin(2\pi h),\qquad Y=s_h+E_s<1.
\]
The supremum is over a complex neighborhood: an error in a nominally real CKKS slot need not be real. Suppose the outer polynomial $p_a$ approximates $\arcsin$ on the real interval $[-s_h,s_h]$ with error $\epsilon_a$, and put
\[
L_a=\sup_{|z|\leq Y}|p_a'(z)|.
\]
If its homomorphic evaluation adds at most $H_a$ before division by $2\pi$, the normalized modular-reduction error is bounded by
\[
E_R\leq\frac{\epsilon_a+L_aE_s+H_a}{2\pi}.
\]
Indeed, first compare the evaluated inner polynomial with $\sin(2\pi t)$; then compare $p_a$ at those two arguments along the complex line segment. The real approximation bound is used only at the ideal real argument. Thus this argument does not silently extend a real minimax approximation guarantee to complex inputs. All circuit errors in $H_s,H_a$ must already be propagated to the relevant output, including rescaling and key switching.

Let $L$ be the ideal final linear transform, including the modular-reduction normalizer and all scale conversions, and let $E_L$ bound its evaluation error. A sufficient base-bootstrap budget is
\[
\|b(c)\|_\infty\leq\varepsilon_{\mathrm{cert}}:=\|L\|_\infty E_R+E_L,
\]
where the right-hand side may be chosen as the certified $\varepsilon$ in the base contract. For example, with an unnormalized coefficient-to-slot evaluation matrix, its norm is at most $n$; a factor $q/\Delta$ in the reconstruction must also appear in $L$. Actual multiplication and rescaling scales obey $\Delta_{\mathrm{mul}}=\Delta_1\Delta_2$ and $\Delta_{\mathrm{rs}}=\Delta_{\mathrm{mul}}/w$ for the actual factor $w$. A scale discrepancy cannot be absorbed into a rounding bound without bounding its message-dependent bias.

These are upper bounds on added error, not proofs of a positive error floor for every input. The error plateau observed for one fixed circuit or parameter set is not a universal lower bound. If $z_0=m+e_{\mathrm{old}}$, the corrected output for $k\geq2$ is $m+e_{\mathrm{old}}-b_k/P_k$; for $k=1$, the added term is $+b_1$ from the base contract. In both cases its total error is at most $\|e_{\mathrm{old}}\|_\infty+\varepsilon/P_k$, and the old error is not recovered. The finite modulus budget $P_kq\leq Q$ also limits this construction. No independence or zero-mean assumption is used. A probabilistic guarantee must control failure on the actual derived inputs, for example through a simultaneous good event or conditional failure bounds given the preceding history; average precision measured on fresh encryptions is insufficient. If the conditional bounds are $\delta_j$, the union bound gives total failure probability at most $\sum_j\delta_j$ without assuming independence. A particular secure parameter set still needs certified domain, arithmetic, and failure-probability bounds; small numerical examples do not supply them.


